\documentclass[11pt,a4paper]{article}
\usepackage[T1]{fontenc}
\usepackage[utf8]{inputenc}
\usepackage[english]{babel}
\usepackage{amsmath,amssymb,amsfonts,amsthm,mathtools}
\usepackage{lmodern,microtype,booktabs,array,longtable}
\usepackage[margin=25mm]{geometry}
\usepackage{xcolor,xurl}
\usepackage[hidelinks,unicode,pdfencoding=auto,bookmarksnumbered]{hyperref}
\usepackage{fancyhdr,etoolbox}
\makeatletter
\patchcmd{\l@section}{1.5em}{2em}{}{}
\makeatother
\setlength{\headheight}{14pt}
\pagestyle{fancy}\fancyhf{}
\fancyhead[L]{\small Supplementary material: gamma condition}
\fancyhead[R]{\small $\alpha_{\rm c}=0.1547$}
\fancyfoot[C]{\thepage}
\setlength{\emergencystretch}{2em}
\allowdisplaybreaks[3]
\raggedbottom
\renewcommand{\thesection}{S\arabic{section}}
\numberwithin{equation}{section}
\numberwithin{table}{section}
\theoremstyle{plain}
\newtheorem{theorem}{Theorem}[section]
\newtheorem{lemma}[theorem]{Lemma}
\newtheorem{proposition}[theorem]{Proposition}
\newtheorem{corollary}[theorem]{Corollary}
\theoremstyle{definition}\newtheorem{definition}[theorem]{Definition}
\theoremstyle{remark}\newtheorem{remark}[theorem]{Remark}
\newcommand{\norm}[1]{\left\|#1\right\|}
\newcommand{\abs}[1]{\left|#1\right|}
\newcommand{\pospart}[1]{\left[#1\right]_+}
\newcommand{\R}{\mathbb R}
\newcommand{\N}{\mathbb N}
\addto\captionsenglish{\renewcommand{\tablename}{Table}\renewcommand{\proofname}{Proof}}
\hypersetup{pdftitle={Supplementary material: semilocal convergence under gamma conditions},
 pdfauthor={Rodrigo Castro Marin; Lia Paz Guzman Pinto; Luis Fernando Echeverri Delgado; Oscar Antonio Restrepo Gutierrez; Gerardo Andres Honorato Gutierrez}}
\title{Supplementary material\\[4pt]\large
Semilocal convergence beyond exact majorization for an efficient three-step method under gamma conditions}
\author{Rodrigo Castro Marín$^1$ \quad Lía Paz Guzmán Pinto$^1$\\
Luis Fernando Echeverri Delgado$^2$\\
Óscar Antonio Restrepo Gutiérrez$^3$ \quad Gerardo Andrés Honorato Gutiérrez$^4$}
\date{}
\begin{document}
\maketitle
\noindent\small
$^1$Institute of Mathematics, Universidad de Valparaíso, Valparaíso, Chile.\\
$^2$Institute of Mathematics, Universidad de Antioquia, Medellín, Colombia.\\
$^3$Institute of Physics, Universidad de Antioquia, Medellín, Colombia.\\
$^4$Institute of Mathematical Engineering, Faculty of Engineering,
Universidad de Valparaíso, Valparaíso, Chile.\\[3pt]
\textbf{Journal:} Numerical Algorithms.\\
\textbf{Corresponding author:} Rodrigo Castro Marín,
\href{mailto:rodrigo.castrom@uv.cl}{rodrigo.castrom@uv.cl}.\\
\normalsize\medskip

This document brings together the algebraic derivations, rational
checks, analysis of previous work, numerical experiments, and generated
data for the semilocal convergence study under gamma conditions.
The programs, data, and execution outputs accompany the document.
Exact comparisons are distinguished from multiprecision floating-point
results.

\setcounter{tocdepth}{2}
\tableofcontents
\clearpage
\section{Derivation of the local error expansion}\label{supp:local-error}
This section derives the local expansion for the three-step method.
The calculation is presented by homogeneous degrees, retaining the
order of the multilinear compositions. Let $X,Y$ be Banach spaces,
$x_*$ a simple solution of $F(x)=0$, and $F\in C^5$ in a neighborhood
of $x_*$. We define
\[
 A=F'(x_*)^{-1},\qquad H_j=\frac1{j!}AF^{(j)}(x_*),\quad j=2,3,4,
 \qquad h_j=H_j(e,\ldots,e),\quad e=x-x_*.
\]
The estimates $O(\|e\|^k)$ are understood in the norm of $X$ or, for
linear operators, in the induced norm. Normalization at the root gives
\begin{equation}
 AF(x_*+e)=e+h_2+h_3+h_4+O(\|e\|^5).
 \label{supp:eq:local-F}
\end{equation}
Likewise,
\[
 AF'(x_*+e)=I+L_1+L_2+L_3+O(\|e\|^4),
\]
\[
 L_1v=2H_2(e,v),\quad L_2v=3H_3(e,e,v),\quad L_3v=4H_4(e,e,e,v).
\]
The Neumann series gives
\begin{equation}
 \begin{aligned}
 B(e):=[AF'(x_*+e)]^{-1}
 ={}&I-L_1+(L_1^2-L_2)\\
 &+(-L_3+L_1L_2+L_2L_1-L_1^3)+O(\|e\|^4).
 \end{aligned}\label{supp:eq:local-inverse}
\end{equation}
We have $F'(x_*+e)^{-1}F(x_*+u)=B(e)AF(x_*+u)$.

\subsection{First stage}
Let $e_y=y-x_*$. Substituting
\eqref{supp:eq:local-F}--\eqref{supp:eq:local-inverse} into the Newton
step,
\begin{equation}
 \begin{aligned}
 e_y={}&h_2+\bigl[2h_3-2H_2(e,h_2)\bigr]\\
 &+\bigl[3h_4-4H_2(e,h_3)+4H_2(e,H_2(e,h_2))-3H_3(e,e,h_2)\bigr]
 +O(\|e\|^5).
 \end{aligned}\label{supp:eq:local-y}
\end{equation}
In particular, $e_y=O(\|e\|^2)$. Since
$AF(y)=e_y+H_2(e_y,e_y)+O(\|e\|^5)$, another application of
\eqref{supp:eq:local-inverse} yields
\begin{equation}
 \begin{aligned}
 B(e)AF(y)={}&h_2+\bigl[2h_3-4H_2(e,h_2)\bigr]\\
 &+\bigl[3h_4-8H_2(e,h_3)+12H_2(e,H_2(e,h_2))\\
 &\hspace{16mm}-6H_3(e,e,h_2)+H_2(h_2,h_2)\bigr]+O(\|e\|^5).
 \end{aligned}\label{supp:eq:local-correction-y}
\end{equation}

\subsection{Second stage}
The second point satisfies $e_z=e_y-5B(e)AF(y)$. By the preceding
identities,
\begin{equation}
 \begin{aligned}
 e_z={}&-4h_2+\bigl[-8h_3+18H_2(e,h_2)\bigr]\\
 &+\bigl[-12h_4+36H_2(e,h_3)-56H_2(e,H_2(e,h_2))\\
 &\hspace{16mm}+27H_3(e,e,h_2)-5H_2(h_2,h_2)\bigr]+O(\|e\|^5).
 \end{aligned}\label{supp:eq:local-z}
\end{equation}
Since $e_z=O(\|e\|^2)$, the only quadratic term contributing up to
degree four in $AF(z)=e_z+H_2(e_z,e_z)+O(\|e\|^5)$ is
$H_2(-4h_2,-4h_2)=16H_2(h_2,h_2)$. Therefore,
\begin{equation}
 \begin{aligned}
 AF(z)={}&-4h_2+\bigl[-8h_3+18H_2(e,h_2)\bigr]\\
 &+\bigl[-12h_4+36H_2(e,h_3)-56H_2(e,H_2(e,h_2))\\
 &\hspace{16mm}+27H_3(e,e,h_2)+11H_2(h_2,h_2)\bigr]+O(\|e\|^5).
 \end{aligned}\label{supp:eq:local-Fz}
\end{equation}

\subsection{Third stage and fourth-degree term}
Relation \eqref{supp:eq:local-y} also gives
$AF(y)=e_y+H_2(h_2,h_2)+O(\|e\|^5)$. Combining this with
\eqref{supp:eq:local-Fz},
\begin{equation}
 \begin{aligned}
 9AF(y)+AF(z)={}&5h_2+10h_3\\
 &+15h_4-20H_2(e,H_2(e,h_2))+20H_2(h_2,h_2)+O(\|e\|^5).
 \end{aligned}\label{supp:eq:local-combination}
\end{equation}
Applying $B(e)/5$ and retaining the terms up to degree four,
\begin{equation}
 \begin{aligned}
 \tfrac15B(e)(9AF(y)+AF(z))={}&h_2+2h_3-2H_2(e,h_2)\\
 &+3h_4-4H_2(e,h_3)-3H_3(e,e,h_2)+4H_2(h_2,h_2)
 +O(\|e\|^5).
 \end{aligned}\label{supp:eq:local-last-correction}
\end{equation}
Subtracting \eqref{supp:eq:local-last-correction} from
\eqref{supp:eq:local-y} cancels all terms of degrees two and three
and gives
\begin{equation}
 e^+=\mathcal Q_4(e)+O(\|e\|^5),\qquad
 \mathcal Q_4(e)=4\bigl[H_2(e,H_2(e,H_2(e,e)))-H_2(H_2(e,e),H_2(e,e))\bigr].
 \label{supp:eq:quartic}
\end{equation}
No associativity of $H_2$ has been assumed. In dimension one,
$H_2(p,q)=c_2pq$, and the two compositions in \eqref{supp:eq:quartic}
coincide; the remainder provides order at least five. In
higher-dimensional spaces, this cancellation is not automatic.

For example, for $F(u,v)=(u+v^2,v+u^2)$, we have
$H_2(p,q)=(p_2q_2,p_1q_1)$ and
$\mathcal Q_4(u,v)=4(v^3-u^3)(u,-v)$. An iteration from $(t,2t)$ gives
\begin{equation}
 T(t,2t)=(28t^4+224t^5,-56t^4+70t^5)+O(t^6).
 \label{supp:eq:path-expansion}
\end{equation}
The program \path{code/verify_audit.py} checks this expansion
and the equivalence of the two expressions for the third stage exactly.

The condition $\mathcal Q_4\not\equiv0$ rules out uniform local order
five for the entire class, but does not determine the order of each
trajectory. If $e_n\neq0$ and
$\liminf\|\mathcal Q_4(e_n)\|/\|e_n\|^4>0$, the remainder bound allows
$O(\|e_n\|^5)$ to be absorbed, yielding
$c\|e_n\|^4\leq\|e_{n+1}\|\leq C\|e_n\|^4$ from some index onward.
The $Q$-order is exactly four when, in addition, the quotient has a
positive finite limit. This distinction is used throughout this document.

\subsection{Exact majorization: consequence for the order and a gamma counterexample}
\label{supp:exact-obstruction}

Two assertions are distinguished. A residual bound along all indices
has consequences for the complete sequence. An expansion of a single
iteration in a direction proves a property of the iteration map, but
does not automatically determine the order of every sequence. The
this distinction is retained throughout the analysis.

\noindent\textbf{Consequence of exact majorization.}
Let $A_0=F'(x_0)^{-1}$, $G=A_0F$, and let $x_*$ be a simple root.
If $x_n\to x_*$, differentiability and the invertibility of $G'(x_*)$
provide constants $c,C>0$ such that
\[
 c\|x_n-x_*\|\leq\|G(x_n)\|\leq C\|x_n-x_*\|
\]
for every sufficiently large index. If $g(t_*)=0$, $g'(t_*)<0$,
$t_n\uparrow t_*$, and $\delta_n=t_*-t_n$, then
$g(t_n)=-g'(t_*)\delta_n+O(\delta_n^2)$.
Therefore, $\|G(x_n)\|\leq g(t_n)$ implies
$\|x_n-x_*\|\leq C_0\delta_n$.

If $\delta_n$ has $Q$-order five, this domination gives $R$-order at
least five. It does not by itself provide $Q$-order five for $x_n$.
The additional bound $\delta_n\leq C_1\|x_n-x_*\|$ would suffice to
obtain the corresponding $Q$-order inequality.

The incompatibility with a fourth-degree lower bound can be seen
without an informal comparison of orders. If $u_{n+1}\geq c u_n^4>0$
and $\delta_{n+1}\leq C\delta_n^5$, where $u_n=\|x_n-x_*\|$, choose
an index $N$ with $c^{1/3}u_N<1$ and $C^{1/4}\delta_N<1$.
For $a=-\log(c^{1/3}u_N)>0$ and
$b=-\log(C^{1/4}\delta_N)>0$, we obtain
\[
 u_{N+j}\geq c^{-1/3}e^{-a4^j},\qquad
 \delta_{N+j}\leq C^{-1/4}e^{-b5^j}.
\]
Hence $u_{N+j}/\delta_{N+j}\to+\infty$. This contradicts domination
by $C_0\delta_n$. The lower-bound assumption must be checked for the
sequence under consideration; it does not follow from $\mathcal Q_4$
being nonzero in some direction.

\medskip\noindent\textbf{A counterexample to universal validity under gamma conditions.}
Consider $F(u,v)=(u+v^2,v+u^2)$, the maximum norm, and $x_0=(h,2h)$,
with $0<h<1/4$. We have
\[
 F'(x_0)=\begin{pmatrix}1&4h\\2h&1\end{pmatrix},\qquad
 A_0=\frac1{1-8h^2}
       \begin{pmatrix}1&-4h\\-2h&1\end{pmatrix}.
\]
From $\|I-F'(x_0)\|_\infty=4h$ and
$\|F(x_0)\|_\infty=2h+h^2$, we obtain the valid choices
\begin{equation}
 \beta_h=\frac{2h+h^2}{1-4h},\qquad
 \gamma_h=\frac1{1-4h},\qquad
 \alpha_h=\frac{2h+h^2}{(1-4h)^2}.
 \label{supp:eq:obstruction-data}
\end{equation}
The map $F''[p,q]=(2p_2q_2,2p_1q_1)$ has norm two and is constant.
Thus $\|A_0F''\|\leq2\gamma_h$ and all its differences are zero.
The gamma condition holds on the required ball, since the domain is
$\mathbb R^2$.

The local expansion obtained in S1 gives
\[
 x_1=(28h^4+224h^5,-56h^4+70h^5)+O(h^6).
\]
Since $A_0=I+O(h)$ and $F(x_1)=x_1+O(h^8)$,
\begin{equation}
 \|A_0F(x_1)\|_\infty=56h^4+O(h^5).
 \label{supp:eq:obstruction-real-residual}
\end{equation}
To compute the scalar residual, we use
$\widetilde g_\alpha(t)=\alpha-t+t^2/(1-t)$ and denote the first
complete iterate from zero by $\widehat t_1$. Substituting the
stages of the method,
\[
 \widehat s_0=\alpha,\qquad
 \widehat v_0=\frac{\alpha(1+4\alpha)}{1-\alpha},\qquad
 \widehat t_1=
 \frac{\alpha(1-2\alpha)(1-2\alpha^2)}
 {(1-\alpha)(1-2\alpha-4\alpha^2)}.
\]
An exact rational simplification yields
\begin{equation}
 \widetilde g_\alpha(\widehat t_1)=
 \frac{2\alpha^5(5-4\alpha-8\alpha^2+8\alpha^3)}
 {(1-\alpha)(1-2\alpha-4\alpha^2)
  (1-4\alpha+6\alpha^3-4\alpha^4)}
 =10\alpha^5+O(\alpha^6).
 \label{supp:eq:obstruction-scalar-residual}
\end{equation}
With \eqref{supp:eq:obstruction-data}, $\alpha_h=2h+O(h^2)$ and
$\gamma_h=1+O(h)$. Since
$\gamma_h g_h(t_1)=\widetilde g_{\alpha_h}(\widehat t_1)$, we obtain
\[
 g_h(t_1)=320h^5+O(h^6),\qquad
 \frac{\|A_0F(x_1)\|_\infty}{g_h(t_1)}
 =\frac{7}{40h}+O(1)\longrightarrow+\infty.
\]
Thus, exact majorization fails at the first iteration for every
sufficiently small $h>0$.

\medskip\noindent\textbf{Rational verification for fixed data.}\par
For $h=1/1000$, the program
\path{code/verify_exact_obstruction.py} computes the stages
using fractions and checks
\[
 \alpha_h=\frac{667}{330672}<\frac{1547}{10000},\qquad
 \|A_0F(x_1)\|_\infty>\frac5{10^{11}},\qquad
 0<g_h(t_1)<\frac4{10^{13}}.
\]
The approximate values, for presentation only, are
\[
 \|A_0F(x_1)\|_\infty\simeq5.599035405559621\,10^{-11},\qquad
 g_h(t_1)\simeq3.367867749356755\,10^{-13}.
\]
The conclusion does not depend on these decimals. The program compares
integer numerators and denominators. The example lies within the
interval of our criterion and refutes universal exact domination by
the proposed majorant; it is not a counterexample to convergence of
the method.



\section{Complete proof for the majorant sequences}\label{supp:scalar}
The algebraic sign proof for the scalar maps is retained. All formulas
in this section correspond to the rational majorant, not to the hatted
controls of the semilocal criterion.
\subsection{Majorant function and normalization}

Let $\beta>0$ and $\gamma>0$. Consider the majorant function 
\begin{equation*}
g(u)=\beta-u+\frac{\gamma u^2}{1-\gamma u}, \qquad 0\leq u<\frac1\gamma.
\end{equation*}
We write 
\begin{equation*}
\alpha=\beta\gamma
\end{equation*}
and assume throughout what follows that 
\begin{equation*}
\alpha<3-2\sqrt2.
\end{equation*}
We denote the small positive root of $g$ by $t_*$.

We also introduce the constant 
\begin{equation*}
\rho=1-\frac1{\sqrt2}.
\end{equation*}

\begin{lemma}[Basic signs of the majorant]\label{supp:s2lemma1}
We have 
\begin{equation*}
0<t_*<\frac{\rho}{\gamma}.
\end{equation*}
Moreover, for every $u\in[0,t_*)$, we have 
\begin{equation*}
g(u)>0,
\end{equation*}
\begin{equation*}
g^{\prime}(u)<0,
\end{equation*}
and 
\begin{equation*}
g^{\prime\prime}(u)>0.
\end{equation*}
\end{lemma}

\begin{proof}
We define the normalized variable 
\begin{equation*}
x=\gamma u.
\end{equation*}
Then 
\begin{equation*}
G(x)=\gamma g\left(\frac{x}{\gamma}\right) =\alpha-x+\frac{x^2}{1-x}.
\end{equation*}
Since 
\begin{equation*}
(1-x)G(x)=\alpha-(1+\alpha)x+2x^2,
\end{equation*}
the assumption 
\begin{equation*}
\alpha<3-2\sqrt2
\end{equation*}
implies that the small positive root $r=\gamma t_*$ satisfies 
\begin{equation*}
0<r<1-\frac1{\sqrt2}=\rho.
\end{equation*}
Therefore, 
\begin{equation*}
0<t_*<\frac{\rho}{\gamma}.
\end{equation*}

Moreover, 
\begin{equation*}
g^{\prime}(u)=-\frac{2\gamma^2u^2-4\gamma u+1}{(1-\gamma u)^2}
\end{equation*}
and 
\begin{equation*}
g^{\prime\prime}(u)=\frac{2\gamma}{(1-\gamma u)^3}>0.
\end{equation*}
If $0\leq u<t_*$, then 
\begin{equation*}
0\leq \gamma u<\gamma t_*=r<\rho.
\end{equation*}
Consequently, 
\begin{equation*}
2\gamma^2u^2-4\gamma u+1>0.
\end{equation*}
It follows that 
\begin{equation*}
g^{\prime}(u)<0.
\end{equation*}
Finally, since $t_*$ is the first positive root of $g$, we have 
\begin{equation*}
g(u)>0, \qquad 0\leq u<t_*.
\end{equation*}
\end{proof}

We define the auxiliary functions 
\begin{equation*}
S(u)=u-\frac{g(u)}{g^{\prime}(u)},
\end{equation*}
\begin{equation*}
V(u)=S(u)-5\frac{g(S(u))}{g^{\prime}(u)},
\end{equation*}
and 
\begin{equation*}
T(u)=S(u)-\frac{9g(S(u))+g(V(u))}{5g^{\prime}(u)}.
\end{equation*}
We also write 
\begin{equation*}
H(u)=g(S(u)), \qquad K(u)=g(V(u)).
\end{equation*}

In the normalized variable 
\begin{equation*}
x=\gamma u, \qquad r=\gamma t_*,
\end{equation*}
we have 
\begin{equation*}
0<r<\rho<\frac12.
\end{equation*}
Since $G(r)=0$, we obtain 
\begin{equation*}
\alpha=r-\frac{r^2}{1-r}=\frac{r(1-2r)}{1-r}.
\end{equation*}
We use the same letters for the normalized functions: 
\begin{equation*}
S(x)=x-\frac{G(x)}{G^{\prime}(x)},
\end{equation*}
\begin{equation*}
V(x)=S(x)-5\frac{G(S(x))}{G^{\prime}(x)},
\end{equation*}
and 
\begin{equation*}
T(x)=S(x)-\frac{9G(S(x))+G(V(x))}{5G^{\prime}(x)}.
\end{equation*}
The relation between the two scales is 
\begin{equation*}
S(\gamma u)=\gamma S(u),
\end{equation*}
\begin{equation*}
V(\gamma u)=\gamma V(u),
\end{equation*}
and 
\begin{equation*}
T(\gamma u)=\gamma T(u),
\end{equation*}
where the normalized functions appear on the left-hand side and the
original functions on the right-hand side.

\subsection{Properties of \texorpdfstring{$S$, $H$ and $V$}{S, H and V}}

\begin{lemma}[Elementary properties of $S$ and $H$]\label{supp:s2lemma2}
For every $u\in[0,t_*)$, we have 
\begin{equation*}
u<S(u)<t_*.
\end{equation*}
Moreover, 
\begin{equation*}
H(u)>0
\end{equation*}
and 
\begin{equation*}
H^{\prime}(u)<0.
\end{equation*}
\end{lemma}

\begin{proof}
Since 
\begin{equation*}
g(u)>0
\end{equation*}
and 
\begin{equation*}
g^{\prime}(u)<0,
\end{equation*}
we have 
\begin{equation*}
S(u)-u=-\frac{g(u)}{g^{\prime}(u)}>0.
\end{equation*}
Therefore, 
\begin{equation*}
u<S(u).
\end{equation*}

On the other hand, 
\begin{equation*}
S(u)=u-\frac{g(u)}{g^{\prime}(u)}.
\end{equation*}
Differentiating, 
\begin{equation*}
S^{\prime}(u)=1-\frac{\left(g^{\prime}(u)\right)^2-g(u)g^{\prime\prime}(u)}{%
\left(g^{\prime}(u)\right)^2}.
\end{equation*}
Hence 
\begin{equation*}
S^{\prime}(u)=\frac{g(u)g^{\prime\prime}(u)}{\left(g^{\prime}(u)\right)^2}.
\end{equation*}
By the preceding lemma, 
\begin{equation*}
g(u)>0,
\end{equation*}
\begin{equation*}
g^{\prime\prime}(u)>0,
\end{equation*}
and 
\begin{equation*}
\left(g^{\prime}(u)\right)^2>0.
\end{equation*}
Thus, 
\begin{equation*}
S^{\prime}(u)>0.
\end{equation*}
The function $S$ is therefore strictly increasing on $[0,t_*]$.
Moreover, 
\begin{equation*}
S(t_*)=t_*-\frac{g(t_*)}{g^{\prime}(t_*)}=t_*.
\end{equation*}
Consequently, if $0\leq u<t_*$, then 
\begin{equation*}
S(u)<S(t_*)=t_*.
\end{equation*}
This proves 
\begin{equation*}
u<S(u)<t_*.
\end{equation*}

Now, 
\begin{equation*}
H(u)=g(S(u)).
\end{equation*}
Since $S(u)<t_*$ and $g>0$ on $[0,t_*)$, we obtain 
\begin{equation*}
H(u)>0.
\end{equation*}
Moreover, 
\begin{equation*}
H^{\prime}(u)=g^{\prime}(S(u))S^{\prime}(u).
\end{equation*}
Since $S(u)<t_*$, we have 
\begin{equation*}
g^{\prime}(S(u))<0.
\end{equation*}
Since also $S^{\prime}(u)>0$, we conclude that 
\begin{equation*}
H^{\prime}(u)<0.
\end{equation*}
\end{proof}

\begin{lemma}[An auxiliary polynomial for the monotonicity of $V$]\label{supp:s2lemma3}
For $0\leq x\leq r$, we have 
\begin{equation*}
R_r(x)>0,
\end{equation*}
where 
\begin{equation*}
\begin{array}{l}
R_r(x)=32r^4x^4-128r^4x^3+216r^4x^2-176r^4x+56r^4 \\ 
\phantom{R_r(x)=}{}-128r^3x^4+428r^3x^3-612r^3x^2+414r^3x-102r^3 \\ 
\phantom{R_r(x)=}{}+204r^2x^4-564r^2x^3+708r^2x^2-426r^2x+87r^2 \\ 
\phantom{R_r(x)=}{}-152rx^4+314rx^3-318rx^2+169rx-31r \\ 
\phantom{R_r(x)=}{}+52x^4-82x^3+60x^2-25x+4.%
\end{array}%
\end{equation*}
\end{lemma}

\begin{proof}
First, 
\begin{equation*}
R_r(r)=4(1-r)^2(2r^2-4r+1)^3>0,
\end{equation*}
because 
\begin{equation*}
2r^2-4r+1>0.
\end{equation*}
Moreover, 
\begin{equation*}
R_r^{\prime}(r)=(r-1)(2r^2-4r+1)^2(32r^2-64r+25).
\end{equation*}
Since 
\begin{equation*}
32r^2-64r+25=16(2r^2-4r+1)+9>0
\end{equation*}
and 
\begin{equation*}
r-1<0,
\end{equation*}
we obtain 
\begin{equation*}
R_r^{\prime}(r)<0.
\end{equation*}

Differentiating three times, 
\begin{equation*}
R_r^{\prime\prime\prime}(x)=12(A(r)x-B(r)),
\end{equation*}
where 
\begin{equation*}
A(r)=64r^4-256r^3+408r^2-304r+104
\end{equation*}
and 
\begin{equation*}
B(r)=64r^4-214r^3+282r^2-157r+41.
\end{equation*}
We have 
\begin{equation*}
A(r)=(2r^2-4r+1)(32r^2-64r+60)+44>0.
\end{equation*}
Moreover, 
\begin{equation*}
B(r)-rA(r)=(1-r)(64r^4-256r^3+366r^2-220r+41).
\end{equation*}
But 
\begin{equation*}
64r^4-256r^3+366r^2-220r+41=(2r^2-4r+1)(32r^2-64r+39)+2>0.
\end{equation*}
Hence 
\begin{equation*}
B(r)-rA(r)>0.
\end{equation*}
If $0\leq x\leq r$, then 
\begin{equation*}
A(r)x-B(r)\leq A(r)r-B(r)<0.
\end{equation*}
Therefore, 
\begin{equation*}
R_r^{\prime\prime\prime}(x)<0, \qquad 0\leq x\leq r.
\end{equation*}
Thus, $R_r^{\prime\prime}$ is decreasing on $[0,r]$.

On the other hand, 
\begin{equation*}
R_r^{\prime\prime}(r)=12(2r^2-4r+1)(16r^4-64r^3+91r^2-54r+10).
\end{equation*}
The second factor satisfies 
\begin{equation*}
16r^4-64r^3+91r^2-54r+10=(2r^2-4r+1)\left(8r^2-16r+\frac{19}{2}%
\right)+\frac12>0.
\end{equation*}
It follows that 
\begin{equation*}
R_r^{\prime\prime}(r)>0.
\end{equation*}
Since $R_r^{\prime\prime}$ is decreasing, for $0\leq x\leq r$ we have 
\begin{equation*}
R_r^{\prime\prime}(x)\geq R_r^{\prime\prime}(r)>0.
\end{equation*}
Hence $R_r^{\prime}$ is increasing. Since $R_r^{\prime}(r)<0$, we obtain 
\begin{equation*}
R_r^{\prime}(x)\leq R_r^{\prime}(r)<0, \qquad 0\leq x\leq r.
\end{equation*}
Therefore, $R_r$ is decreasing on $[0,r]$. Finally, 
\begin{equation*}
R_r(x)\geq R_r(r)>0, \qquad 0\leq x\leq r.
\end{equation*}
\end{proof}

\begin{lemma}[Monotonicity of $V$]\label{supp:s2lemma4}
The function $V$ is strictly decreasing on $[0,r)$. In particular, 
\begin{equation*}
V^{\prime}(x)<0, \qquad 0\leq x<r.
\end{equation*}
\end{lemma}

\begin{proof}
A direct calculation gives 
\begin{equation*}
V^{\prime}(x)=\frac{2(r-x)(1-2r-2x+2rx)R_r(x)}{(r-1)(2x^2-4x+1)^3D_r(x)^2},
\end{equation*}
where 
\begin{equation*}
D_r(x)=2r^2x-2r^2-4rx+2r+3x-1.
\end{equation*}
That is, 
\begin{equation*}
D_r(x)=x(2r^2-4r+3)-(2r^2-2r+1).
\end{equation*}
Since 
\begin{equation*}
2r^2-4r+3>0
\end{equation*}
and 
\begin{equation*}
x\leq r,
\end{equation*}
we have 
\begin{equation*}
D_r(x)\leq D_r(r).
\end{equation*}
Now, 
\begin{equation*}
D_r(r)=(r-1)(2r^2-4r+1)<0.
\end{equation*}
Hence 
\begin{equation*}
D_r(x)<0,
\end{equation*}
and, in particular, 
\begin{equation*}
D_r(x)^2>0.
\end{equation*}

For $0\leq x<r$, we have 
\begin{equation*}
r-x>0.
\end{equation*}
Moreover, 
\begin{equation*}
1-2r-2x+2rx\geq1-4r+2r^2=2r^2-4r+1>0.
\end{equation*}
By the preceding lemma, 
\begin{equation*}
R_r(x)>0.
\end{equation*}
Also, 
\begin{equation*}
r-1<0,
\end{equation*}
\begin{equation*}
2x^2-4x+1>0,
\end{equation*}
and 
\begin{equation*}
D_r(x)^2>0.
\end{equation*}
Therefore, the numerator of $V^{\prime}(x)$ is positive and its
denominator is negative. It follows that 
\begin{equation*}
V^{\prime}(x)<0, \qquad 0\leq x<r.
\end{equation*}
\end{proof}

\begin{lemma}[Location of $V$]\label{supp:s2lemma5}
For every $0\leq x<r$, we have 
\begin{equation*}
r<V(x)<1.
\end{equation*}
Equivalently, for every $0\leq u<t_*$, 
\begin{equation*}
t_*<V(u)<\frac1\gamma.
\end{equation*}
\end{lemma}

\begin{proof}
Since 
\begin{equation*}
G(r)=0,
\end{equation*}%
we have 
\begin{equation*}
S(r)=r.
\end{equation*}
Hence 
\begin{equation*}
V(r)=S(r)-5\frac{G(S(r))}{G^{\prime }(r)}=r.
\end{equation*}%
By the preceding lemma, $V$ is decreasing. Therefore, if $0\leq x<r$,
then 
\begin{equation*}
V(x)>V(r)=r.
\end{equation*}

It remains to prove that $V(x)<1$. Since $V$ is decreasing and $x\geq0$,
it suffices to prove 
\begin{equation*}
V(0)<1.
\end{equation*}
A direct calculation gives 
\begin{equation*}
1-V(0)=\frac{16r^4-12r^3-3r^2+4r-1}{(r-1)(2r^2-2r+1)}.
\end{equation*}
Let 
\begin{equation*}
N(r)=16r^4-12r^3-3r^2+4r-1.
\end{equation*}
Then 
\begin{equation*}
N^{\prime}(r)=64r^3-36r^2-6r+4.
\end{equation*}
Factoring, 
\begin{equation*}
N^{\prime}(r)=2(2r-1)(16r^2-r-2).
\end{equation*}
Since 
\begin{equation*}
0<r<\rho<\frac12,
\end{equation*}
we have 
\begin{equation*}
2r-1<0.
\end{equation*}
Moreover, 
\begin{equation*}
16r^2-r-2<16\rho^2-2=22-16\sqrt2<0.
\end{equation*}
Therefore, 
\begin{equation*}
N^{\prime}(r)>0.
\end{equation*}
Thus, $N$ is increasing on $(0,\rho)$. Hence 
\begin{equation*}
N(r)<N(\rho).
\end{equation*}
Now, 
\begin{equation*}
N(\rho)=\frac{73-52\sqrt2}{2}<0.
\end{equation*}
It follows that 
\begin{equation*}
N(r)<0.
\end{equation*}
Since 
\begin{equation*}
r-1<0
\end{equation*}
and 
\begin{equation*}
2r^2-2r+1>0,
\end{equation*}
the denominator of $1-V(0)$ is negative. Consequently, 
\begin{equation*}
1-V(0)>0.
\end{equation*}
This proves that 
\begin{equation*}
V(0)<1.
\end{equation*}
Since $V(x)\leq V(0)$, we obtain 
\begin{equation*}
V(x)<1.
\end{equation*}
Therefore, 
\begin{equation*}
r<V(x)<1.
\end{equation*}
Dividing by $\gamma>0$ gives 
\begin{equation*}
t_*<V(u)<\frac1\gamma.
\end{equation*}
\end{proof}

\subsection{Properties of \texorpdfstring{$T$}{T}}

\begin{lemma}[An inequality for $K=g\circ V$]\label{supp:s2lemma6}
For every $u\in[0,t_*)$, we have 
\begin{equation*}
K(u)>-4H(u).
\end{equation*}
Consequently, 
\begin{equation*}
9H(u)+K(u)>0.
\end{equation*}
\end{lemma}

\begin{proof}
Since 
\begin{equation*}
V(u)<\frac1\gamma,
\end{equation*}
the quantity $g(V(u))$ is well defined. By the convexity of $g$, 
\begin{equation*}
g(V(u))\geq g(S(u))+g^{\prime}(S(u))(V(u)-S(u)).
\end{equation*}
That is, 
\begin{equation*}
K(u)\geq H(u)+g^{\prime}(S(u))(V(u)-S(u)).
\end{equation*}
But 
\begin{equation*}
V(u)-S(u)=-5\frac{g(S(u))}{g^{\prime}(u)}=-5\frac{H(u)}{g^{\prime}(u)}.
\end{equation*}
Substituting, 
\begin{equation*}
K(u)\geq H(u)-5g^{\prime}(S(u))\frac{H(u)}{g^{\prime}(u)}.
\end{equation*}
Therefore, 
\begin{equation*}
K(u)\geq H(u)\left(1-5\frac{g^{\prime}(S(u))}{g^{\prime}(u)}\right).
\end{equation*}
Since 
\begin{equation*}
u<S(u)<t_*,
\end{equation*}
and $g^{\prime\prime}>0$, the function $g^{\prime}$ is strictly
increasing. Thus, 
\begin{equation*}
g^{\prime}(u)<g^{\prime}(S(u))<0.
\end{equation*}
Dividing by $g^{\prime}(u)<0$ gives 
\begin{equation*}
0<\frac{g^{\prime}(S(u))}{g^{\prime}(u)}<1.
\end{equation*}
Hence 
\begin{equation*}
1-5\frac{g^{\prime}(S(u))}{g^{\prime}(u)}>-4.
\end{equation*}
Since $H(u)>0$, we conclude that 
\begin{equation*}
K(u)>-4H(u).
\end{equation*}
Consequently, 
\begin{equation*}
9H(u)+K(u)>9H(u)-4H(u)=5H(u)>0.
\end{equation*}
\end{proof}

\begin{lemma}[The inequality $S<T$]\label{supp:s2lemma7}
For every $u\in[0,t_*)$, we have 
\begin{equation*}
S(u)<T(u).
\end{equation*}
\end{lemma}

\begin{proof}
By definition, 
\begin{equation*}
T(u)-S(u)=-\frac{9g(S(u))+g(V(u))}{5g^{\prime}(u)}.
\end{equation*}
Using 
\begin{equation*}
H(u)=g(S(u))
\end{equation*}
and 
\begin{equation*}
K(u)=g(V(u)),
\end{equation*}
we obtain 
\begin{equation*}
T(u)-S(u)=-\frac{9H(u)+K(u)}{5g^{\prime}(u)}.
\end{equation*}
By the preceding lemma, 
\begin{equation*}
9H(u)+K(u)>0.
\end{equation*}
Moreover, 
\begin{equation*}
g^{\prime}(u)<0.
\end{equation*}
Therefore, 
\begin{equation*}
T(u)-S(u)>0.
\end{equation*}
That is, 
\begin{equation*}
S(u)<T(u).
\end{equation*}
\end{proof}

\begin{lemma}[Sign of $Q_r$]\label{supp:s2lemma8}
For every $0\leq x\leq r$, we have 
\begin{equation*}
Q_r(x)>0,
\end{equation*}
where 
\begin{equation*}
\begin{array}{l}
Q_r(x)=8r^4x^4-32r^4x^3+24r^4x^2+16r^4x-16r^4 \\ 
\phantom{Q_r(x)=}{}-32r^3x^4+152r^3x^3-168r^3x^2+36r^3x+12r^3 \\ 
\phantom{Q_r(x)=}{}+36r^2x^4-216r^2x^3+252r^2x^2-84r^2x+3r^2 \\ 
\phantom{Q_r(x)=}{}-8rx^4+116rx^3-132rx^2+46rx-4r \\ 
\phantom{Q_r(x)=}{}-2x^4-28x^3+30x^2-10x+1.%
\end{array}%
\end{equation*}
\end{lemma}

\begin{proof}
Evaluating at $x=r$ gives 
\begin{equation*}
Q_r(r)=(1-r)^2(2r^2-4r+1)^3>0.
\end{equation*}
We make the change 
\begin{equation*}
y=r-x.
\end{equation*}
Then 
\begin{equation*}
Q_r^{\prime}(r-y)=-2\widehat H_r(y),
\end{equation*}
where 
\begin{equation*}
\widehat H_r(y)=4A(r)y^3+3B(r)y^2+2C(r)y+E(r),
\end{equation*}
with 
\begin{equation*}
A(r)=4r^4-16r^3+18r^2-4r-1,
\end{equation*}
\begin{equation*}
B(r)=-16r^5+80r^4-148r^3+124r^2-54r+14,
\end{equation*}
\begin{equation*}
C(r)=24r^6-144r^5+348r^4-432r^3+294r^2-108r+15,
\end{equation*}
and 
\begin{equation*}
E(r)=-16r^7+112r^6-324r^5+500r^4-440r^3+216r^2-53r+5.
\end{equation*}

First we prove that 
\begin{equation*}
A(r)<0.
\end{equation*}
Indeed, 
\begin{equation*}
A^{\prime}(r)=4(r-1)(4r^2-8r+1).
\end{equation*}
The function $A$ decreases and then increases on $[0,\rho]$.
Moreover, 
\begin{equation*}
A(0)=-1
\end{equation*}
and 
\begin{equation*}
A(\rho)=-1.
\end{equation*}
Consequently, 
\begin{equation*}
A(r)\leq -1<0.
\end{equation*}

Now, 
\begin{equation*}
\widehat H_r^{\prime}(y)=12A(r)y^2+6B(r)y+2C(r).
\end{equation*}
Since 
\begin{equation*}
A(r)<0,
\end{equation*}
the function $\widehat H_r^{\prime}$ is concave in $y$. Moreover, 
\begin{equation*}
\widehat H_r^{\prime}(0)=2C(r).
\end{equation*}
Factoring, 
\begin{equation*}
\widehat H_r^{\prime}(0)=6(2r^2-6r+5)(2r^2-4r+1)(2r^2-2r+1)>0.
\end{equation*}
Also, 
\begin{equation*}
\widehat H_r^{\prime}(r)=6(4r^4-28r^3+42r^2-22r+5)>0.
\end{equation*}
To justify the last inequality, we define 
\begin{equation*}
M(r)=4r^4-28r^3+42r^2-22r+5.
\end{equation*}
Then 
\begin{equation*}
M(\rho)=-7+6\sqrt2>0.
\end{equation*}
Moreover, 
\begin{equation*}
M^{\prime}(r)=16r^3-84r^2+84r-22
\end{equation*}
and 
\begin{equation*}
M^{\prime\prime}(r)=48r^2-168r+84>0
\end{equation*}
on $[0,\rho]$. Thus, $M^{\prime}$ is increasing on $[0,\rho]$,
and therefore 
\begin{equation*}
M^{\prime}(r)\leq M^{\prime}(\rho)=-24+14\sqrt2<0.
\end{equation*}
Hence $M$ is decreasing and, since $r<\rho$, we obtain 
\begin{equation*}
M(r)>M(\rho)>0.
\end{equation*}
It follows that 
\begin{equation*}
\widehat H_r^{\prime}(r)>0.
\end{equation*}

Since $\widehat H_r^{\prime}$ is concave and positive at the endpoints
of $[0,r]$, we obtain 
\begin{equation*}
\widehat H_r^{\prime}(y)>0, \qquad 0\leq y\leq r.
\end{equation*}
Hence $\widehat H_r$ is increasing. Moreover, 
\begin{equation*}
\widehat H_r(0)=E(r).
\end{equation*}
But 
\begin{equation*}
E(r)=-(r-1)(2r^2-4r+1)^2(4r^2-8r+5)>0.
\end{equation*}
Therefore, 
\begin{equation*}
\widehat H_r(y)>0, \qquad 0\leq y\leq r.
\end{equation*}
Thus, 
\begin{equation*}
Q_r^{\prime}(r-y)=-2\widehat H_r(y)<0.
\end{equation*}
Hence $Q_r$ is decreasing with respect to $x$. If $0\leq x<r$, then 
\begin{equation*}
Q_r(x)>Q_r(r)>0.
\end{equation*}
For $x=r$, we already know that $Q_r(r)>0$. This completes the proof.
\end{proof}

\begin{lemma}[Sign of $P_{r}$]\label{supp:s2lemma9}
For every $0\leq x\leq r$, we have 
\begin{equation*}
P_{r}(x)<0,
\end{equation*}%
where 
\begin{equation*}
\begin{array}{l}
P_{r}(x)=8r^{3}x^{5}-40r^{3}x^{4}+56r^{3}x^{3}-8r^{3}x^{2}-32r^{3}x+16r^{3}
\\ 
\phantom{P_r(x)=}{}%
+24r^{2}x^{6}-168r^{2}x^{5}+424r^{2}x^{4}-424r^{2}x^{3}+100r^{2}x^{2}+72r^{2}x-28r^{2}
\\ 
\phantom{P_r(x)=}{}-48rx^{6}+276rx^{5}-548rx^{4}+382rx^{3}-2rx^{2}-79rx+19r
\\ 
\phantom{P_r(x)=}{}+20x^{6}-92x^{5}+122x^{4}-6x^{3}-63x^{2}+33x-5.%
\end{array}%
\end{equation*}
\end{lemma}

\begin{proof}
We write 
\begin{equation*}
a=1-r, \qquad z=1-x.
\end{equation*}
Then 
\begin{equation*}
\frac1{\sqrt2}<a<1, \qquad a\leq z\leq 1.
\end{equation*}
In these variables, 
\begin{equation*}
P_r(x)=\Pi(a,z),
\end{equation*}
where 
\begin{equation*}
\begin{array}{l}
\Pi(a,z)=8a^3z^5-24a^3z^3+24a^2z^6-56a^2z^4+52a^2z^2 \\ 
\phantom{\Pi(a,z)=}{}-36az^5+62az^3-37az-4z^6+18z^4-21z^2+9.%
\end{array}%
\end{equation*}
Evaluating at $z=a$, we obtain 
\begin{equation*}
\Pi(a,a)=(2a-3)(2a+3)(2a^2-1)^3<0.
\end{equation*}
Therefore, it suffices to prove that $\Pi(a,z)$ is strictly decreasing
with respect to $z$ on the interval $[a,1]$.

Differentiating, 
\begin{equation*}
\begin{array}{l}
\dfrac{\partial \Pi}{\partial z}%
(a,z)=40a^3z^4-72a^3z^2+144a^2z^5-224a^2z^3+104a^2z \\ 
\phantom{\dfrac{\partial \Pi}{\partial z}(a,z)=}{}%
-180az^4+186az^2-37a-24z^5+72z^3-42z.%
\end{array}%
\end{equation*}
We prove that 
\begin{equation*}
-\frac{\partial \Pi}{\partial z}(a,z)>0.
\end{equation*}
Since $a\leq z$, we set 
\begin{equation*}
y=\frac{a}{z}.
\end{equation*}
Then 
\begin{equation*}
0<y\leq1,
\end{equation*}
and, since $a>1/\sqrt2$, we have 
\begin{equation*}
y>\frac1{\sqrt2 z}.
\end{equation*}
Moreover, $1/\sqrt2<z\leq1$. A direct calculation gives 
\begin{equation*}
-\frac{\partial \Pi}{\partial z}(a,z)=z\mathcal{L}(y,z),
\end{equation*}
where 
\begin{equation*}
\begin{array}{l}
\mathcal{L}(y,z)=-40y^3z^6-144y^2z^6+72y^3z^4+224y^2z^4+180yz^4+24z^4 \\ 
\phantom{\mathcal L(y,z)=}{}-104y^2z^2-186yz^2-72z^2+37y+42.%
\end{array}%
\end{equation*}
Since $z>0$, it suffices to prove that 
\begin{equation*}
\mathcal{L}(y,z)>0.
\end{equation*}

Differentiating with respect to $y$, 
\begin{equation*}
\begin{array}{l}
\dfrac{\partial \mathcal{L}}{\partial y}%
(y,z)=-120z^6y^2-288z^6y+216z^4y^2+448z^4y \\ 
\phantom{\dfrac{\partial \mathcal L}{\partial y}(y,z)=}{}%
+180z^4-208z^2y-186z^2+37.%
\end{array}%
\end{equation*}
Moreover, 
\begin{equation*}
\frac{\partial^2 \mathcal{L}}{\partial y^2}(y,z)=-16z^2%
\left(15z^4y+18z^4-27z^2y-28z^2+13\right).
\end{equation*}
Let 
\begin{equation*}
B_z(y)=15z^4y+18z^4-27z^2y-28z^2+13.
\end{equation*}
The coefficient of $y$ in $B_z$ is 
\begin{equation*}
z^2(15z^2-27)<0.
\end{equation*}
Therefore, $B_z$ is decreasing in $y$. Since 
\begin{equation*}
y\geq \frac1{\sqrt2 z},
\end{equation*}
we have 
\begin{equation*}
B_z(y)\leq B_z\left(\frac1{\sqrt2 z}\right).
\end{equation*}
Now, 
\begin{equation*}
B_z\left(\frac1{\sqrt2 z}\right)=18z^4+\frac{15}{\sqrt2}z^3-28z^2-\frac{27}{%
\sqrt2}z+13.
\end{equation*}
Denote 
\begin{equation*}
\Phi(z)=18z^4+\frac{15}{\sqrt2}z^3-28z^2-\frac{27}{\sqrt2}z+13.
\end{equation*}
Then 
\begin{equation*}
\Phi^{\prime\prime}(z)=216z^2+45\sqrt2 z-56>0
\end{equation*}
for $1/\sqrt2\leq z\leq1$. Consequently, $\Phi$ is convex.
Moreover, 
\begin{equation*}
\Phi\left(\frac1{\sqrt2}\right)=-\frac{25}{4}<0
\end{equation*}
and 
\begin{equation*}
\Phi(1)=3-6\sqrt2<0.
\end{equation*}
Being convex, $\Phi$ attains its maximum at the endpoints of the
interval. Hence 
\begin{equation*}
\Phi(z)<0.
\end{equation*}
Therefore, 
\begin{equation*}
B_z(y)<0.
\end{equation*}
It follows that 
\begin{equation*}
\frac{\partial^2 \mathcal{L}}{\partial y^2}(y,z)>0.
\end{equation*}
Thus, the function $\partial \mathcal{L}/\partial y$ is increasing
with respect to $y$. Consequently, 
\begin{equation*}
\frac{\partial \mathcal{L}}{\partial y}(y,z)\geq \frac{\partial \mathcal{L}}{%
\partial y}\left(\frac1{\sqrt2 z},z\right).
\end{equation*}
A direct calculation gives 
\begin{equation*}
\begin{array}{l}
\dfrac{\partial \mathcal{L}}{\partial y}\left(\dfrac1{\sqrt2
z},z\right)=-144\sqrt2 z^5+120z^4+224\sqrt2 z^3 \\ 
\phantom{\dfrac{\partial \mathcal L}{\partial y}\left(\dfrac1{\sqrt2
z},z\right)=}{}-78z^2-104\sqrt2 z+37.%
\end{array}%
\end{equation*}
Moreover, 
\begin{equation*}
\begin{array}{l}
-144\sqrt2 z^5+120z^4+224\sqrt2 z^3-78z^2-104\sqrt2 z+37 \\ 
=-144\sqrt2\left(z-\frac1{\sqrt2}\right)\left(z^4+\frac{\sqrt2}{12}z^3-\frac{%
53}{36}z^2-\frac{67\sqrt2}{144}z+\frac{37}{144}\right).%
\end{array}%
\end{equation*}
For $1/\sqrt2\leq z\leq1$, we have 
\begin{equation*}
\begin{array}{l}
z^4+\dfrac{\sqrt2}{12}z^3-\dfrac{53}{36}z^2-\dfrac{67\sqrt2}{144}z+\dfrac{37%
}{144} \\ 
\leq z^2+\dfrac{\sqrt2}{12}z-\dfrac{53}{36}z^2-\dfrac{67\sqrt2}{144}z+\dfrac{%
37}{144} \\ 
=-\dfrac{17}{36}z^2-\dfrac{55\sqrt2}{144}z+\dfrac{37}{144} \\ 
\leq -\dfrac{17}{72}-\dfrac{55}{144}+\dfrac{37}{144}<0.%
\end{array}%
\end{equation*}
Therefore, 
\begin{equation*}
\frac{\partial \mathcal{L}}{\partial y}\left(\frac1{\sqrt2 z},z\right)\geq0.
\end{equation*}
Hence 
\begin{equation*}
\frac{\partial \mathcal{L}}{\partial y}(y,z)>0.
\end{equation*}
Thus, $\mathcal{L}$ is increasing with respect to $y$. Consequently, 
\begin{equation*}
\mathcal{L}(y,z)\geq \mathcal{L}\left(\frac1{\sqrt2 z},z\right).
\end{equation*}
A direct calculation gives 
\begin{equation*}
\mathcal{L}\left(\frac1{\sqrt2 z},z\right)=-48z^4+80\sqrt2
z^3+40z^2-75\sqrt2 z-10+\frac{37\sqrt2}{2z}.
\end{equation*}
We define 
\begin{equation*}
\Lambda(z)=-48z^4+80\sqrt2 z^3+40z^2-75\sqrt2 z-10+\frac{37\sqrt2}{2z}.
\end{equation*}
Then 
\begin{equation*}
\Lambda\left(\frac1{\sqrt2}\right)=0.
\end{equation*}
Moreover, 
\begin{equation*}
\begin{array}{l}
2z^2\Lambda^{\prime}(z)=-384z^5+480\sqrt2 z^4+160z^3-150\sqrt2 z^2-37\sqrt2
\\ 
=-384\left(z-\frac1{\sqrt2}\right)\left(z^4-\frac{3\sqrt2}{4}z^3-\frac76z^2-%
\frac{37\sqrt2}{192}z-\frac{37}{192}\right).%
\end{array}%
\end{equation*}
Since 
\begin{equation*}
z^4-\frac76z^2\leq -\frac16z^2<0,
\end{equation*}
we have 
\begin{equation*}
z^4-\frac{3\sqrt2}{4}z^3-\frac76z^2-\frac{37\sqrt2}{192}z-\frac{37}{192}<0.
\end{equation*}
Therefore, 
\begin{equation*}
\Lambda^{\prime}(z)\geq0
\end{equation*}
for $1/\sqrt2\leq z\leq1$. Since $z>1/\sqrt2$, we conclude that 
\begin{equation*}
\Lambda(z)>\Lambda\left(\frac1{\sqrt2}\right)=0.
\end{equation*}
Thus, 
\begin{equation*}
\mathcal{L}(y,z)>0.
\end{equation*}
Consequently, 
\begin{equation*}
-\frac{\partial \Pi}{\partial z}(a,z)>0.
\end{equation*}
Hence 
\begin{equation*}
\frac{\partial \Pi}{\partial z}(a,z)<0.
\end{equation*}
Therefore, $\Pi(a,z)$ is strictly decreasing with respect to $z$.
Since $z\geq a$, we obtain 
\begin{equation*}
\Pi(a,z)\leq\Pi(a,a)<0.
\end{equation*}
Thus, 
\begin{equation*}
P_r(x)<0.
\end{equation*}
\end{proof}

\begin{lemma}[The inequality $T(x)<r$]\label{supp:s2lemma10}
For every $0\leq x<r$, we have 
\begin{equation*}
T(x)<r.
\end{equation*}
Equivalently, for every $0\leq u<t_*$, 
\begin{equation*}
T(u)<t_*.
\end{equation*}
\end{lemma}

\begin{proof}
A direct algebraic calculation gives 
\begin{equation*}
T(x)-r=\frac{2(r-x)^5P_r(x)}{(r-1)(2x^2-4x+1)^3D_r(x)Q_r(x)},
\end{equation*}
where 
\begin{equation*}
D_r(x)=2r^2x-2r^2-4rx+2r+3x-1.
\end{equation*}
We have already proved that 
\begin{equation*}
D_r(x)<0,
\end{equation*}
\begin{equation*}
Q_r(x)>0,
\end{equation*}
and 
\begin{equation*}
P_r(x)<0.
\end{equation*}
Moreover, for $0\leq x<r$, 
\begin{equation*}
(r-x)^5>0.
\end{equation*}
Therefore, the numerator of $T(x)-r$ is negative.

On the other hand, 
\begin{equation*}
r-1<0,
\end{equation*}
\begin{equation*}
2x^2-4x+1>0,
\end{equation*}
\begin{equation*}
D_r(x)<0,
\end{equation*}
and 
\begin{equation*}
Q_r(x)>0.
\end{equation*}
Hence the denominator is positive. Consequently, 
\begin{equation*}
T(x)-r<0.
\end{equation*}
That is, 
\begin{equation*}
T(x)<r.
\end{equation*}
Dividing by $\gamma>0$ gives 
\begin{equation*}
T(u)<t_*.
\end{equation*}
\end{proof}

\subsection{Functional proposition and scalar sequence}

\begin{proposition}[Functional properties]\label{supp:s2functional}
For every $u\in[0,t_*)$, the following properties hold: 
\begin{equation*}
u<S(u)<t_*,
\end{equation*}
\begin{equation*}
H(u)>0, \qquad H^{\prime}(u)<0,
\end{equation*}
\begin{equation*}
t_*<V(u)<\frac1\gamma,
\end{equation*}
\begin{equation*}
K(u)>-4H(u),
\end{equation*}
\begin{equation*}
9H(u)+K(u)>0,
\end{equation*}
\begin{equation*}
S(u)<T(u)<t_*,
\end{equation*}
and 
\begin{equation*}
V^{\prime}(u)<0.
\end{equation*}
In particular, $V$ is strictly decreasing on $[0,t_*)$.
\end{proposition}

\begin{proof}
The first two properties follow from Lemma~\ref{supp:s2lemma2}.
The location of $V$ follows from Lemma~\ref{supp:s2lemma5}.
The inequalities 
\begin{equation*}
K(u)>-4H(u)
\end{equation*}
and 
\begin{equation*}
9H(u)+K(u)>0
\end{equation*}
follow from Lemma~\ref{supp:s2lemma6}. The inequality 
\begin{equation*}
S(u)<T(u)
\end{equation*}
follows from Lemma~\ref{supp:s2lemma7}, while 
\begin{equation*}
T(u)<t_*
\end{equation*}
follows from Lemma~\ref{supp:s2lemma10}. Finally, the strict
monotonicity of $V$ follows from Lemma~\ref{supp:s2lemma4}, transferring
the inequality from the normalized variable $x$ to the original
variable $u$.
\end{proof}

We define 
\begin{equation*}
t_0=0,
\end{equation*}
\begin{equation*}
s_n=S(t_n),
\end{equation*}
\begin{equation*}
v_n=V(t_n),
\end{equation*}
and 
\begin{equation*}
t_{n+1}=T(t_n).
\end{equation*}
Equivalently, 
\begin{equation*}
s_n=t_n-\frac{g(t_n)}{g^{\prime}(t_n)},
\end{equation*}
\begin{equation*}
v_n=s_n-5\frac{g(s_n)}{g^{\prime}(t_n)},
\end{equation*}
and 
\begin{equation*}
t_{n+1}=s_n-\frac{9g(s_n)+g(v_n)}{5g^{\prime}(t_n)}.
\end{equation*}

\begin{theorem}[Interlacing of the sequences]\label{supp:s2chain}
The preceding sequence is well defined and, for every $n\geq0$,
satisfies 
\begin{equation*}
0\leq t_n\leq s_n\leq t_{n+1}\leq s_{n+1}<t_*<v_{n+1}\leq v_n.
\end{equation*}
Moreover, 
\begin{equation*}
t_*<v_n<\frac1\gamma, \qquad n\geq0.
\end{equation*}
\end{theorem}

\begin{proof}
We proceed by induction.

Since 
\begin{equation*}
t_0=0<t_*,
\end{equation*}
we can apply the preceding proposition with $u=t_0$. We obtain 
\begin{equation*}
t_0<S(t_0)<T(t_0)<t_*.
\end{equation*}
Since 
\begin{equation*}
s_0=S(t_0)
\end{equation*}
and 
\begin{equation*}
t_1=T(t_0),
\end{equation*}
we have 
\begin{equation*}
0\leq t_0\leq s_0\leq t_1<t_*.
\end{equation*}
Applying the proposition again, now with $u=t_1$, gives 
\begin{equation*}
t_1\leq s_1<t_*.
\end{equation*}
We also obtain 
\begin{equation*}
t_*<v_1<\frac1\gamma.
\end{equation*}
On the other hand, the proposition applied at $u=t_0$ gives 
\begin{equation*}
t_*<v_0<\frac1\gamma.
\end{equation*}
Since $V$ is decreasing and 
\begin{equation*}
t_0\leq t_1,
\end{equation*}
we deduce 
\begin{equation*}
v_1=V(t_1)\leq V(t_0)=v_0.
\end{equation*}
Therefore, 
\begin{equation*}
0\leq t_0\leq s_0\leq t_1\leq s_1<t_*<v_1\leq v_0.
\end{equation*}
This proves the base case.

Now suppose that, for some $n\geq0$, we have 
\begin{equation*}
0\leq t_n<t_*.
\end{equation*}
Applying the proposition with $u=t_n$, we obtain 
\begin{equation*}
t_n<S(t_n)<T(t_n)<t_*.
\end{equation*}
Since 
\begin{equation*}
s_n=S(t_n)
\end{equation*}
and 
\begin{equation*}
t_{n+1}=T(t_n),
\end{equation*}
we obtain 
\begin{equation*}
0\leq t_n\leq s_n\leq t_{n+1}<t_*.
\end{equation*}
In particular, 
\begin{equation*}
0\leq t_{n+1}<t_*.
\end{equation*}
Applying the proposition again, now with $u=t_{n+1}$, gives 
\begin{equation*}
t_{n+1}\leq S(t_{n+1})<t_*.
\end{equation*}
Since 
\begin{equation*}
s_{n+1}=S(t_{n+1}),
\end{equation*}
we have 
\begin{equation*}
t_{n+1}\leq s_{n+1}<t_*.
\end{equation*}
Moreover, by the property of $V$, 
\begin{equation*}
t_*<V(t_{n+1})<\frac1\gamma.
\end{equation*}
That is, 
\begin{equation*}
t_*<v_{n+1}<\frac1\gamma.
\end{equation*}
Finally, since $V$ is decreasing and 
\begin{equation*}
t_n\leq t_{n+1},
\end{equation*}
it follows that 
\begin{equation*}
v_{n+1}=V(t_{n+1})\leq V(t_n)=v_n.
\end{equation*}
Collecting the preceding inequalities, we conclude that 
\begin{equation*}
0\leq t_n\leq s_n\leq t_{n+1}\leq s_{n+1}<t_*<v_{n+1}\leq v_n.
\end{equation*}
Moreover, at each step we have 
\begin{equation*}
v_n<\frac1\gamma,
\end{equation*}
so all terms $g(v_n)$ appearing in the definition of $t_{n+1}$ are
well defined. This completes the induction.
\end{proof}


\subsection{Convergence and asymptotic constants}
For the following references, we write the chain already proved as
\begin{equation}
0\leq t_n\leq s_n\leq t_{n+1}\leq s_{n+1}<t_*<v_{n+1}\leq v_n<1/\gamma.
\label{supp:eq:scalar-chain}
\end{equation}
\begin{theorem}[Convergence and asymptotic constants]\label{supp:s2asymptotics}
The sequences in Theorem~\ref{supp:s2chain} satisfy $t_n\uparrow t_*$,
$s_n\uparrow t_*$, and $v_n\downarrow t_*$. Let $r=\gamma t_*$
and let $d=2r^2-4r+1>0$. Then
\begin{equation}
\begin{aligned}
\lim_{n\to\infty}\frac{t_*-s_n}{(t_*-t_n)^2}
 &=\frac{\gamma}{(1-r)d},\qquad
\lim_{n\to\infty}\frac{v_n-t_*}{(t_*-t_n)^2}
 =\frac{4\gamma}{(1-r)d},\\
\lim_{n\to\infty}\frac{t_*-t_{n+1}}{(t_*-t_n)^5}
 &=\frac{2\gamma^4(1+2r)(5-2r)}{(1-r)^4d^4}.
\end{aligned}
\label{supp:eq:scalar-asymptotic-constants}
\end{equation}
In particular, there exist $C_s,C_v,C_5>0$ and $n_0\geq0$ such that,
for $n\geq n_0$,
\begin{equation}
\begin{aligned}
0\leq t_*-s_n&\leq C_s(t_*-t_n)^2,\qquad
0\leq v_n-t_*\leq C_v(t_*-t_n)^2,\\
0\leq t_*-t_{n+1}&\leq C_5(t_*-t_n)^5.
\end{aligned}
\label{supp:eq:scalar-error-bounds}
\end{equation}
\end{theorem}
\begin{proof}
\textit{Convergence of the three sequences.}
The chain \eqref{supp:eq:scalar-chain} shows that $\{t_n\}$ and
$\{s_n\}$ are increasing and bounded above by $t_*$, while $\{v_n\}$
is decreasing and bounded below by $t_*$. Let
$\ell=\lim_{n\to\infty}t_n\leq t_*$. The maps $S,V,T$ extend
continuously to $t_*$ because $g'(t_*)\neq0$, and satisfy
$S(t_*)=V(t_*)=T(t_*)=t_*$. Taking the limit in $t_{n+1}=T(t_n)$
gives $\ell=T(\ell)$. If $\ell<t_*$,
Proposition~\ref{supp:s2functional} would give $T(\ell)>\ell$,
a contradiction. Therefore, $\ell=t_*$. Continuity now gives
\[
 s_n=S(t_n)\longrightarrow S(t_*)=t_*,\qquad
 v_n=V(t_n)\longrightarrow V(t_*)=t_*.
\]

\textit{Asymptotic constants.}
To distinguish the scales explicitly in this calculation, we write
$\widetilde g=G$ and
\[
 \widehat S(x)=\gamma S(x/\gamma),\quad
 \widehat V(x)=\gamma V(x/\gamma),\quad
 \widehat T(x)=\gamma T(x/\gamma),
\]
where the right-hand sides contain the original functions of $g$.
The hatted functions are the normalized maps used in the preceding
factorizations. With $x=\gamma t$ and $r=\gamma t_*$, recall that
\[
 \widetilde g'(r)=-\frac{d}{(1-r)^2},\qquad
 \widetilde g''(r)=\frac{2}{(1-r)^3},\qquad d=2r^2-4r+1>0.
\]
Set $\delta=r-x$ and $a=1/((1-r)d)$. The expansion of the Newton step
at the simple zero $r$ gives
\[
 r-\widehat S(x)
 =-\frac{\widetilde g''(r)}{2\widetilde g'(r)}\delta^2
   +O(\delta^3)
 =a\delta^2+O(\delta^3).
\]
In particular, $\widehat S(x)-r=O(\delta^2)$. Expanding $\widetilde g$
at $r$ and using $\widetilde g'(x)=\widetilde g'(r)+O(\delta)$ yields
\[
 \frac{\widetilde g(\widehat S(x))}{\widetilde g'(x)}
 =\widehat S(x)-r+O(\delta^3)
 =-a\delta^2+O(\delta^3).
\]
Substituting into the second stage,
\[
 \widehat V(x)-r
 =\widehat S(x)-r
   -5\frac{\widetilde g(\widehat S(x))}{\widetilde g'(x)}
 =4a\delta^2+O(\delta^3).
\]
For the third stage, the expansion at $x=r$ of the exact factorization
of $\widehat T(x)-r$ used in the preceding proof gives
\[
 r-\widehat T(x)
 =\frac{2(1+2r)(5-2r)}{(1-r)^4d^4}\delta^5+O(\delta^6).
\]
The identities yielding this constant are
\[
 D_r(r)=(r-1)d,\qquad Q_r(r)=(1-r)^2d^3,\qquad
 P_r(r)=(2r-5)(2r+1)d^3.
\]
They are checked by substituting $x=r$ into the preceding polynomials.
The program \path{code/symbolic_majorant_verification.py} verifies
these three identities and the complete factorization of $T(x)-r$.

To return to the original sequences, we take $x=\gamma t_n$, divide
each expansion by $\gamma$, and substitute $\delta=\gamma(t_*-t_n)$.
The constants of the quadratic terms are multiplied by $\gamma$, and
that of the fifth-degree term by $\gamma^4$. This yields the three
limits in \eqref{supp:eq:scalar-asymptotic-constants}. All are positive
and finite, since $0<r<1-1/\sqrt2$. Finally, each quotient is bounded
by a positive constant from some index onward; taking a common index
gives the three estimates in \eqref{supp:eq:scalar-error-bounds}.
\end{proof}
\subsection{Summability of the scalar controls}

To transfer the scalar control to the problem in the Banach space, we
will need the lengths of the three stages and the residual combination
appearing in the last step. We define
\begin{equation*}
\begin{aligned}
a_n&=s_n-t_n,\qquad b_n=v_n-s_n,\qquad
c_n=t_{n+1}-s_n,\\
D_n&=9g(s_n)+g(v_n).
\end{aligned}
\end{equation*}

\begin{corollary}[Summability of the scalar controls]
\label{supp:cor:global-sums}
For every $p>0$,
\[
\sum_{n=0}^{\infty}(t_*-t_n)^p<\infty.
\]
Moreover,
\begin{equation}
\begin{aligned}
&\sum_{n=0}^{\infty}a_n\leq t_*,\qquad
\sum_{n=0}^{\infty}c_n\leq t_*,\qquad
\sum_{n=0}^{\infty}b_n<\infty,
\qquad \sum_{n=0}^{\infty}b_n^2<\infty,\\
&\sum_{n=0}^{\infty}g(t_n)\leq t_*,\qquad
\sum_{n=0}^{\infty}D_n\leq5t_*.
\end{aligned}
\label{supp:eq:global-sums}
\end{equation}
\end{corollary}

\begin{proof}
Let $e_n=t_*-t_n$. By \eqref{supp:eq:scalar-error-bounds}, there exists
$n_1$ such that $e_{n+1}\leq C_5e_n^5$ and $C_5e_n^4\leq1/2$ for
$n\geq n_1$. Hence $e_{n+1}\leq e_n/2$, and geometric comparison
gives the convergence of $\sum e_n^p$ for every $p>0$.

Interlacing gives $0\leq a_n,c_n\leq t_{n+1}-t_n$; on summing, the
first two bounds in \eqref{supp:eq:global-sums} are telescopic.
Moreover,
\[
b_n=(v_n-t_*)+(t_*-s_n)=O(e_n^2),
\]
and the summability of $e_n^2$ and $e_n^4$ gives the assertions about
$b_n$. Finally,
\[
g(t_n)=(-g'(t_n))a_n\leq a_n,
\qquad
D_n=5(-g'(t_n))c_n\leq5c_n,
\]
because $0<-g'(t_n)\leq1$. Summing gives the two remaining bounds.
\end{proof}



\section{Radius controls, first cycle, and quantitative comparison}
\label{supp:radial}
This section develops the rational verification of the criterion
$0<\alpha=\beta\gamma\leq1547/10000$ and the quantitative comparisons.
The rational recurrences $\tau_n,\varepsilon_n$ are distinguished
from the sequences $\widehat\tau_n,\widehat\varepsilon_n$, which use
a specific estimate for the first cycle. Both retain the same scalar
majorant $g$ and the same method; the hats do not denote a change
of algorithm.

\subsection{Assumptions and estimates in the normalized scale}
Let $G(x)=F'(x_0)^{-1}F(x)$,
$\rho=1-1/\sqrt2$, $R_\gamma=\rho/\gamma$, and
$g(t)=\beta-t+\gamma t^2/(1-\gamma t)$. We assume
$B[x_0,R_\gamma]\subset\Omega$, $\|G(x_0)\|\leq\beta$,
$\|G''(x_0)\|\leq2\gamma$, and, for
$x,z\in B[x_0,R_\gamma]$ and $u,v\in[0,R_\gamma]$ with
$\|x-z\|\leq|u-v|$,
\[
 \|G''(x)-G''(z)\|\leq|g''(u)-g''(v)|.
\]
Set $\mathcal F(\xi)=\gamma G(x_0+\xi/\gamma)$. Then
$\mathcal F'(0)=I$, $\|\mathcal F(0)\|\leq\alpha=\beta\gamma$,
and $\|\mathcal F''(0)\|\leq2$. The difference condition above implies
\begin{equation}
 \|\mathcal F''(\xi)-\mathcal F''(\zeta)\|\leq6\|\xi-\zeta\|,
 \qquad \xi,\zeta\in B[0,\rho].
 \label{supp:eq:normalized-Lipschitz}
\end{equation}
Indeed, subdividing the segment between the points into $N$ parts and
applying the gamma condition with parameters $d/N$ and $0$ gives
$N(M(d/N)-M(0))\to6d$, where $d=\|\xi-\zeta\|$ and
$M(t)=2/(1-t)^3$. For large $N$, the parameters belong to the required
domain. The converse uses $M'(t)\geq6$.
Only $F\in C^2$ is assumed, not the existence of $F'''$.

The rational functions used in the cycles are
\begin{equation}
 \begin{aligned}
 k(t)&=\frac{t^2}{1-t},&m(t)&=2-\frac1{(1-t)^2},&M(t)&=\frac2{(1-t)^3},\\
 L(t,d)&=\frac{d(2(1-t)-d)}{(1-t)^2(1-t-d)^2},&
 E_2(t,d)&=\frac{d^2}{(1-t)^2(1-t-d)},&
 E_3(t,d)&=\frac{d^3}{(1-t)^3(1-t-d)}.
 \end{aligned}\label{supp:eq:rational-controls}
\end{equation}
For $0\leq t<\rho$, $m(t)>0$. If $\|\xi\|\leq t$, then
$\|\mathcal F'(\xi)^{-1}\|\leq1/m(t)$.
The functions $L,E_2,E_3$ control, respectively, the variation of the
derivative, the first-order remainder, and the remainder after
subtracting the quadratic term. Their integral representations are
\[
 L(t,d)=\int_0^dM(t+s)\,ds,\quad
 E_2(t,d)=\int_0^d(d-s)M(t+s)\,ds,
\]
\[
 E_3(t,d)=\int_0^d(d-s)(M(t+s)-M(t))\,ds.
\]
They apply when $t+d\leq\rho$ and give
\begin{equation}
\begin{aligned}
 \|\mathcal F'(\xi+h)-\mathcal F'(\xi)\|&\leq L(t,d),\\
 \|\mathcal F(\xi+h)-\mathcal F(\xi)-\mathcal F'(\xi)h\|&\leq E_2(t,d),\\
 \|\mathcal F(\xi+h)-\mathcal F(\xi)-\mathcal F'(\xi)h
 -\tfrac12\mathcal F''(\xi)[h,h]\|&\leq E_3(t,d),
\end{aligned}
\label{supp:eq:remainder-bounds}
\end{equation}
whenever $\|\xi\|\leq t$ and $\|h\|\leq d$. The monotonicity and
convexity of $M$ prove that $L,E_2,E_3$ are nondecreasing in each
argument. When $t+d\leq\sigma<\rho$, we have
$L(t,d)\leq M(\sigma)d$ and $E_2(t,d)\leq M(\sigma)d^2/2$.

\subsection{Residual identities and rational recurrences}
For a cycle, we write $J=\mathcal F'(\xi)$, $p=y-\xi$, $q=z-y$,
$r=\xi^+-y$, $w=r-q/5$, $D=\mathcal F'(y)-J$, and
\[
 \mathcal B(d)=\mathcal F(y+d)-\mathcal F(y)-\mathcal F'(y)d,
 \qquad \mathcal E(d)=\mathcal B(d)-\tfrac12\mathcal F''(y)[d,d].
\]
The relations $Jq=-5\mathcal F(y)$ and
$Jr=-(9\mathcal F(y)+\mathcal F(z))/5$ give the identities
\begin{equation}
 \begin{aligned}
 w&=-\tfrac15J^{-1}(Dq+\mathcal B(q)),\\
 \mathcal F(\xi^+)&=Dw+\mathcal B(r)-\tfrac15\mathcal B(q)\\
 &=Dw+\tfrac12\mathcal F''(y)[w,w]+\tfrac15\mathcal F''(y)[w,q]
   -\tfrac2{25}\mathcal F''(y)[q,q]+\mathcal E(r)-\tfrac15\mathcal E(q).
 \end{aligned}\label{supp:eq:residual-identities}
\end{equation}
The last equality uses $r=w+q/5$ and bilinear symmetry, without
reassociating compositions of $\mathcal F''$. First, $y,z$ are placed
in the domain; the first identity controls $\xi^+$ before its
residual is evaluated.

Given controls $(\tau_n,\varepsilon_n)$, the rational recurrences are
\begin{equation}
 \begin{aligned}
 m_n&=m(\tau_n),&A_n&=\varepsilon_n/m_n,&S_n&=\tau_n+A_n,\\
 Y_n&=E_2(\tau_n,A_n),&B_n&=5Y_n/m_n,&V_n&=S_n+B_n,\\
 \ell_n&=L(\tau_n,A_n),&W_n&=\frac{\ell_nB_n+E_2(S_n,B_n)}{5m_n},&
 C_n&=B_n/5+W_n,\\
 \tau_{n+1}&=S_n+C_n.&&&&
 \end{aligned}\label{supp:eq:radial-recursion}
\end{equation}
We define
\begin{equation}
 \begin{aligned}
 P_n^{(2)}&=\ell_nW_n+E_2(S_n,C_n)+\tfrac15E_2(S_n,B_n),\\
 P_n^{(3)}&=\ell_nW_n+M(S_n)\left(\tfrac12W_n^2+
  \tfrac15W_nB_n+\tfrac2{25}B_n^2\right)
  +E_3(S_n,C_n)+\tfrac15E_3(S_n,B_n),\\
 \varepsilon_{n+1}&=\min\{P_n^{(2)},P_n^{(3)}\}.
 \end{aligned}\label{supp:eq:radial-residual}
\end{equation}
A cycle is admissible if its quantities are constructed in the indicated
order, $\varepsilon_n\geq0$, and $\tau_n,S_n,V_n,\tau_{n+1}<\rho$.
For an admissible cycle with
$\|\xi_n\|\leq\tau_n$ and
$\|\mathcal F(\xi_n)\|\leq\varepsilon_n$, the inverse bound,
\eqref{supp:eq:remainder-bounds}, and
\eqref{supp:eq:residual-identities} give
\begin{equation}
\begin{aligned}
 \|\bar p_n\|&\leq A_n,& \|\upsilon_n\|&\leq S_n,&
 \|\mathcal F(\upsilon_n)\|&\leq Y_n,\\
 \|\bar q_n\|&\leq B_n,& \|\zeta_n\|&\leq V_n,&
 \|\bar r_n-\bar q_n/5\|&\leq W_n,\\
 \|\bar r_n\|&\leq C_n,& \|\xi_{n+1}\|&\leq\tau_{n+1},&
 \|\mathcal F(\xi_{n+1})\|&\leq\varepsilon_{n+1}.
\end{aligned}
\label{supp:eq:one-step-bounds}
\end{equation}
Monotonicity of the controls follows successively from that of
$1/m,L,E_2,E_3$ and the nonnegativity of all products; the minimum of
two nondecreasing quantities is also nondecreasing. Thus, the upper
data of an admissible cycle dominate those of any cycle started below
them.

\subsection{Specific estimate for the first cycle}
We fix
\[
 \alpha_{\rm c}=\frac{1547}{10000},\qquad
 \sigma_{\rm all}=\frac{3661}{12500},\qquad
 \sigma_{\rm tail}=\frac6{25}.
\]

For the initial cycle, we use the sharper consequence of
\eqref{supp:eq:normalized-Lipschitz}, namely
$\|\mathcal F''(\xi)\|\leq2+6\|\xi\|$. By integration, the
three remainder estimates in \eqref{supp:eq:remainder-bounds} also
hold, for $t+d\leq\rho$, with the controls
\begin{equation}
\begin{aligned}
 M_*(t)&=2+6t,& L_*(t,d)&=(2+6t)d+3d^2,\\
 E_{2,*}(t,d)&=(1+3t)d^2+d^3,& E_{3,*}(d)&=d^3.
\end{aligned}
\label{supp:eq:first-polynomial-controls}
\end{equation}
Indeed, we integrate $2+6t+6\theta d$ and
$(1-\theta)(2+6t+6\theta d)$, respectively; for the second-order
remainder we integrate $6\theta(1-\theta)d^3$. No $F'''$ is required.

We use hats to distinguish the sequences that employ these bounds only
at $n=0$. We define
\begin{equation}
\begin{aligned}
 \widehat\tau_0&=0,&
 \widehat\varepsilon_0=\widehat A_0=\widehat S_0&=\alpha,&
 \widehat m_0&=1,\\
 \widehat Y_0&=\alpha^2+\alpha^3,&
 \widehat B_0&=5\widehat Y_0,&
 \widehat V_0&=\alpha+\widehat B_0,\\
 \widehat\ell_0&=2\alpha+3\alpha^2,&
 \widehat W_0&=\frac{\widehat\ell_0\widehat B_0+
 E_{2,*}(\alpha,\widehat B_0)}5,&
 \widehat C_0&=\widehat Y_0+\widehat W_0,\\
 \widehat\tau_1&=\alpha+\widehat C_0.&&&
\end{aligned}
\label{supp:eq:first-cycle-controls}
\end{equation}
The next residual is controlled by
\begin{equation}
\begin{aligned}
 \widehat\varepsilon_1={}&\widehat\ell_0\widehat W_0
 +(2+6\alpha)\left(\frac{\widehat W_0^2}{2}
 +\frac{\widehat W_0\widehat B_0}{5}
 +\frac{2\widehat B_0^2}{25}\right)
 +\widehat C_0^3+\frac{\widehat B_0^3}{5}.
\end{aligned}
\label{supp:eq:first-cycle-residual}
\end{equation}
For $n\geq1$, we use exactly
\eqref{supp:eq:radial-recursion}--\eqref{supp:eq:radial-residual},
placing a hat on every quantity in the cycle. The functions
$m,M,L,E_2,E_3$ and the iteration of $F$ do not change.

\begin{lemma}[Control of the first cycle]\label{supp:lem:first-cycle-controls}
If $F$ satisfies the gamma condition and $0<\alpha\leq\alpha_{\rm c}$,
the initial stages are well defined and satisfy
\eqref{supp:eq:one-step-bounds} with the controls in
\eqref{supp:eq:first-cycle-controls}--\eqref{supp:eq:first-cycle-residual}.
In particular,
\begin{equation}
 \widehat V_0\leq0.292871921615<\sigma_{\rm all}<\rho,\qquad
 \widehat\tau_1<\frac15,\qquad
 \widehat\varepsilon_1<\frac1{75}.
\label{supp:eq:first-cycle-upper-bounds}
\end{equation}
\end{lemma}
\begin{proof}
Since $\mathcal F'(0)=I$, $\|\upsilon_0\|\leq\alpha$, and the
Newton cancellation gives
\[
 \|\mathcal F(\upsilon_0)\|
 \leq\alpha^2\int_0^1(1-\theta)(2+6\theta\alpha)\,d\theta
 =\alpha^2+\alpha^3=\widehat Y_0.
\]
Then $\|\bar q_0\|\leq\widehat B_0$ and
$\|\zeta_0\|\leq\widehat V_0=\alpha+5\alpha^2+5\alpha^3$.
This polynomial is increasing, and its value at $\alpha_{\rm c}$ is
$0.292871921615$. The inequality $\sigma_{\rm all}<\rho$ is checked
exactly through $2(1-\sigma_{\rm all})^2-1=2921/78125000>0$.
Thus, $\mathcal F(\zeta_0)$ can be evaluated within the domain.

The first identity in \eqref{supp:eq:residual-identities}, together
with \eqref{supp:eq:first-polynomial-controls}, gives
$\|\bar r_0-\bar q_0/5\|\leq\widehat W_0$ and
$\|\bar r_0\|\leq\widehat C_0$. Thus,
$\|\xi_1\|\leq\widehat\tau_1$ before its residual is evaluated.
All quantities in
\eqref{supp:eq:first-cycle-controls}--\eqref{supp:eq:first-cycle-residual}
are polynomials in $\alpha$ with nonnegative coefficients. Their
rational evaluation at $\alpha_{\rm c}$ gives
\[
 \widehat\tau_1<0.198986457<1/5,\qquad
 \widehat\varepsilon_1<0.013187894<1/75.
\]
The first bound also places the segment from $\upsilon_0$ to $\xi_1$
inside the ball. Applying \eqref{supp:eq:residual-identities}, with
the remainders bounded by $\widehat C_0^3$ and $\widehat B_0^3$,
then gives $\|\mathcal F(\xi_1)\|\leq\widehat\varepsilon_1$.
The fractions are recorded in the results file described below.
\end{proof}

We also define
\begin{equation}
\begin{aligned}
 \widehat\Phi_{\rm ef}(\alpha)&=\sup_{n\geq0}
  \max\{\widehat V_n,\widehat\tau_{n+1}\},\\
 \widehat\Delta_n&=\frac{[\widehat\varepsilon_n-\gamma g(t_n)]_+}{\gamma},&
 \widehat E_n&=\frac{[\widehat A_n+\widehat C_n-
 \gamma(t_{n+1}-t_n)]_+}{\gamma}.
\end{aligned}
\label{supp:eq:first-cycle-defects}
\end{equation}
The same admissibility convention is adopted for
$\widehat\Phi_{\rm ef}$. For $0<\alpha\leq\alpha_{\rm c}$,
\eqref{supp:eq:one-step-bounds} holds at $n=0$ by
Lemma~\ref{supp:lem:first-cycle-controls} and for $n\geq1$ by the
standard cycle estimates above. Moreover,
$\widehat\tau_{n+1}-\widehat\tau_n=\widehat A_n+\widehat C_n$ and
$\widehat\varepsilon_n\leq\widehat A_n$ also hold at $n=0$.
Hence boundedness of the hatted radii gives telescopic summability of
$\widehat A_n+\widehat C_n$ and of the corresponding nonnegative
defects. The one-step bounds then make $\{x_n\}$ a Cauchy sequence;
continuity gives a zero of $F$, and the standard averaged-derivative
argument with $k'(t)<1$ gives uniqueness in the gamma ball. The same
telescoping argument yields the a priori error bounds, while the
residual-distance estimate gives the residual error bound. These are
the properties used in the finite tail verification below, which
starts at an index $N\geq1$.

\subsection{Finite verification and control of all remaining indices}
The tail control is based on the uniform constants $m=m(\sigma)$,
$M=M(\sigma)$ and the polynomials
\begin{equation}
 \begin{aligned}
 a(e)&=e/m,&b(e)&=\frac{5M}{2m}a(e)^2,\\
 w(e)&=\frac{Ma(e)b(e)+Mb(e)^2/2}{5m},&c(e)&=b(e)/5+w(e),\\
 P(e)&=Ma(e)w(e)+\frac M2\bigl(c(e)^2+b(e)^2/5\bigr).
 \end{aligned}\label{supp:eq:tail-controls}
\end{equation}
We have
\begin{equation}
 \begin{aligned}
 P(e)={}&\frac{5M^3}{4m^6}e^4+\frac{7M^4}{8m^8}e^5+
 \frac{7M^5}{16m^{10}}e^6+\frac{5M^6}{16m^{12}}e^7+
 \frac{25M^7}{128m^{14}}e^8.
 \end{aligned}\label{supp:eq:tail-polynomial}
\end{equation}
The program \path{code/verify_algebra.py} checks this
identity symbolically.

If $\widehat\tau_N\leq T$ and $\widehat\varepsilon_N\leq e$, it
suffices to find $q\in(0,1)$ and $\sigma<\rho$ such that
\begin{equation}
 P(e)\leq qe,\qquad
 T+\max\{(a(e)+c(e))/(1-q),a(e)+b(e)\}<\sigma.
 \label{supp:eq:tail-test}
\end{equation}
To explain why all indices are controlled, set $D=a(e)+c(e)$,
$H=a(e)+b(e)$. The nonnegative coefficients give
$a(q^je)+c(q^je)\leq q^jD$, $a(q^je)+b(q^je)\leq q^jH$, and
$P(q^je)\leq q^jP(e)$. Assuming inductively that
\[
 \widehat\tau_{N+j}\leq T+D\frac{1-q^j}{1-q},\qquad
 \widehat\varepsilon_{N+j}\leq q^je,
\]
we obtain
\[
 \widehat\tau_{N+j}+a(q^je)+b(q^je)
 \leq T+\frac D{1-q}(1-q^j)+Hq^j<\sigma.
\]
This places the first two stages before their remainders are estimated.
Then $\widehat W_{N+j}\leq w(q^je)$ and
$\widehat C_{N+j}\leq c(q^je)$ allow the third stage to be placed and
the residual to be bounded by $P(q^je)\leq q^{j+1}e$. Induction gives
\begin{equation}
 \widehat\tau_\infty\leq T+\frac D{1-q},\qquad
 \sum_{i=j}^{\infty}(\widehat A_{N+i}+\widehat C_{N+i})
 \leq\frac{Dq^j}{1-q}.
 \label{supp:eq:tail-consequence}
\end{equation}
This is the uniform estimate, not an extrapolation from a finite sample.

\subsection{Rational states and value of the criterion}
The polynomials of the first cycle are evaluated at $\alpha_{\rm c}$.
All the following comparisons are between exact fractions. We obtain
\[
 \widehat V_0=0.292871921615,\quad
 \widehat\tau_1<0.198986457<1/5,\quad
 \widehat\varepsilon_1<0.013187894<1/75.
\]
We therefore start the upper sequence of the cycles $n=1,2$ at
$(\overline\tau_1,\overline\varepsilon_1)=(1/5,1/75)$.
At the end of each cycle, $\mathcal R(u)=10^{-24}\lceil10^{24}u\rceil$
is applied to the data passed to the next cycle. The exact transmitted
states are
\begin{equation}
 \begin{aligned}
 \overline\tau_2&=\frac{236509841095667281330063}{10^{24}},&
 \overline\varepsilon_2&=\frac{426195191308547945528}{10^{24}},\\
 \overline\tau_3&=\frac{238026135303284715478357}{10^{24}},&
 \overline\varepsilon_3&=\frac{5818357024425655}{10^{24}}.
 \end{aligned}\label{supp:eq:exact-states}
\end{equation}
The numerators and denominators of all intermediate quantities, not
only these four states, are included in
\path{results/criterion_alpha1547.json}.

We take $T=238026136/10^9$, $e=6/10^9$, $\sigma=6/25$, and $q=1/2$.
In this case, $m(\sigma)=97/361$ and $M(\sigma)=31250/6859$.
The computed fractions satisfy
\begin{equation}
 \begin{aligned}
 P(e)/e&<6.8\times10^{-20}<1/2,\\
 T+2(a(e)+c(e))&<0.238026181<0.24,\\
 T+a(e)+b(e)&<0.238026159<0.24.
 \end{aligned}\label{supp:eq:tail-values}
\end{equation}
We also verify that
\[
 (3-\alpha_{\rm c})^2-8=\frac{9573209}{10^8}>0,\qquad
 2(1-\sigma_{\rm all})^2-1=\frac{2921}{78125000}>0.
\]
All radii of the cycles $n=1,2$ are less than $0.253$.
Together with the first cycle and \eqref{supp:eq:tail-values}, this
gives $\widehat\Phi_{\rm ef}(\alpha_{\rm c})\leq\sigma_{\rm all}<\rho$.
Monotonicity extends the result to $0<\alpha\leq\alpha_{\rm c}$.
The hatted version of the semilocal theorem gives, in particular,
\[
 \|x_*-x_0\|<\frac{0.238026181}{\gamma},\qquad
 x_n,y_n,z_n\in B[x_0,0.24/\gamma]\quad(n\geq3).
\]
For $n=3+j$, we also have the bound
$\|x_*-x_{3+j}\|\leq2(a(e)+c(e))2^{-j}/\gamma$.

\subsection{Quantitative comparison with Cordero et al. (2015)}
We retain the functions of \cite{CorderoSemilocal2015}:
\[
 h_{15}(a)=a/2+a^2/2+5a^3/8,\quad
 f_{15}(a)=[1-a(1+h_{15}(a))]^{-1},
\]
\[
 g_{15}(a)=a/2+(1+a)h_{15}(a)+ah_{15}(a)^2/2.
\]
The first positive zero of $f_{15}(a)^2g_{15}(a)-1$ on the admissible
branch satisfies
\[
 a_{15}^*\in(0.29313992158369,0.29313992158370).
\]
The endpoints are rational; the signs are checked exactly. On the
positive branch, $h_{15},f_{15},g_{15}$ are increasing, and
$f_{15}^2g_{15}$ is strictly increasing; therefore, the sign change
identifies the endpoint of the interval starting from zero.

For $F_\alpha(u,v)=(u+v^2-\alpha,v+u^2-\alpha/2)$, with the maximum
norm and $x_0=0$, the gamma condition uses $\beta=\alpha$, $\gamma=1$.
The smallest constants for the 2015 criterion are $\beta_{15}=1$,
$\eta_{15}=\alpha$, $K=2$, so $a=2\alpha$.
The admitted interval is $0<\alpha<a_{15}^*/2\simeq0.146569960792$,
strictly contained in $0<\alpha\leq0.1547$.
For $\alpha=3/20$, $a=3/10$ and
$f_{15}(3/10)^2g_{15}(3/10)>1.06695>1$.
Failure of a sufficient condition does not imply divergence of the method.

For the location comparison, we use the unhatted rational controls,
started at $(0,\alpha)$, as indicated in the table in this section.
Their limit $\tau_\infty$ is not identified with $\widehat\tau_\infty$.
The program \path{code/compare_cordero2015.py} performs a finite
verification for each datum and applies \eqref{supp:eq:tail-test}
to that sequence. The 2015 bounds for the limit and for all stages are
\[
 L_{15}(a)=\frac{(1+h_{15}(a))\alpha}{1-f_{15}(a)g_{15}(a)},\qquad
 R_{15}(a)\alpha=\left(\frac{5a}{2}+
 \frac{1+h_{15}(a)}{1-f_{15}(a)g_{15}(a)}\right)\alpha.
\]
The recurrence-based sum for the location is
$\sum_{n\geq0}(1+h_{15}(a_n))d_n$, with $a_0=a$, $d_0=\alpha$,
$a_{n+1}=a_nf_{15}(a_n)^2g_{15}(a_n)$, and
$d_{n+1}=f_{15}(a_n)g_{15}(a_n)d_n$.
Its first two terms sum to
\[
 Q_{15}^{[2]}=(1+h_{15}(a))\alpha+
 (1+h_{15}(a_1))f_{15}(a)g_{15}(a)\alpha.
\]
In all common tabulated cases, fractions are used to check
$U_\alpha<Q_{15}^{[2]}<L_{15}<R_{15}\alpha$, where
$\tau_\infty\leq U_\alpha$. Thus, the improvement does not come from
comparing only with a larger geometric majorization.

\begin{table}[htbp]
\centering\small
\begin{tabular}{@{}lccc@{}}\toprule
$\alpha$ & Upper bound for $\tau_\infty$ & $Q_{15}^{[2]}$ & $L_{15}$\\\midrule
$0.01$ & $0.010103200647$ & $0.010310404439$ & $0.010316962789$\\
$0.02$ & $0.020427400093$ & $0.021286724646$ & $0.021344357094$\\
$0.05$ & $0.053036118831$ & $0.059049322277$ & $0.060285803938$\\
$0.1$ & $0.116956870642$ & $0.146445008466$ & $0.166567164179$\\
$0.14$ & $0.195079729913$ & $0.262956274107$ & $0.406176058866$\\
$0.15$ & $0.230706675248$ & -- & --\\
\bottomrule\end{tabular}
\caption{The bound in the first numerical column is rounded upward;
the other two columns are approximations. The signs are checked using
complete fractions in the program and the JSON file, not these rounded
values. The 2015 criterion does not apply at $\alpha=0.15$.}
\label{supp:tab:localization}
\end{table}

\subsection{Programs and scope of the verification}
\path{code/verify_gamma_criterion.py} verifies the first cycle,
the rational states of the subsequent cycles, and all tail inequalities.
\path{code/compare_cordero2015.py} checks the separation of the
criteria and the location bounds. Both verifications use
\texttt{fractions.Fraction}; they do not decide signs through
floating-point arithmetic. Decimals are generated after the comparisons.
The justification for all indices is the tail induction, not the
execution of a finite number of steps.


\section{Audit of the previous works from 2015, 2016, and 2022}\label{supp:audit}
We examine the published results of Cordero et
al.~\cite{CorderoSemilocal2015}, Singh et al.\ \cite{SinghEtAl2016},
and Maroju, Behl, and Motsa \cite{MarojuBehlMotsa2022}.
The cited pages are the printed pages of those works. The aim is to
clarify which assertions can be used and which corrections their
proofs require. No intentions are attributed to the authors, and the
falsity of an entire theorem is not inferred from any typographical error.

Four situations are distinguished: an identity or inequality refuted
by exact calculation; an insufficient argument; a false conclusion
with a counterexample satisfying the assumptions; and a numerical
value that does not agree with the published formula. The programs
separate rational and symbolic checks from multiprecision reproductions;
the latter are not interval-arithmetic results.

\subsection{Equivalence of the method and obstruction to fifth order}
In the three articles, the last stage is written as
\[
 x^+=z-\tfrac15\Gamma(-16F(y)+F(z)),\qquad
 z=y-5\Gamma F(y),\qquad \Gamma=F'(x)^{-1}.
\]
Substituting $z$ gives $x^+=y-\Gamma(9F(y)+F(z))/5$, exactly the scheme
studied here. The different expressions do not produce different
methods. The expansion in Section~\ref{supp:local-error} contains
\[
 \mathcal Q_4(e)=4\{H_2(e,H_2(e,H_2(e,e)))-H_2(H_2(e,e),H_2(e,e))\}.
\]
For $F(u,v)=(u+v^2,v+u^2)$, $H_2(p,q)=(p_2q_2,p_1q_1)$ and
$\mathcal Q_4(u,v)=4(v^3-u^3)(u,-v)$.
By \eqref{supp:eq:path-expansion},
$\|T(t,2t)\|_\infty/\|(t,2t)\|_\infty^4\to7/2$, whereas the
quotient with the fifth power diverges. The root is simple and the
operator is polynomial. Uniform local fifth order cannot be assigned
to the general class. This does not determine a single order for all
trajectories or by itself refute semilocal convergence.

\subsection{Cordero et al. (2015): corrections and valid semilocal content}
\paragraph{A1. The function $h$ retains the factor of the second stage.}
On pages 206--207 of \cite{CorderoSemilocal2015}, the definition is
\[
 h(a)=a/2+a^2/2+5a^3/8=\frac a{10}\bigl[4+(1+5a/2)^2\bigr].
\]
This identity is exact. For $F_a(u)=u-1-au^2/2$ and $x_0=0$, we
obtain $y_0=1$, $z_0=1+5a/2$, and $x_1=1+h(a)$. The bound is exact
in that example. The missing factor identified in the 2022 work
should not be attributed to the 2015 work.

\paragraph{A2. An intermediate inequality in (6) is false.}
The chain on page 207 includes
\[
 \|\Gamma_0\|\,\|F(x_0)+(9F(y_0)+F(z_0))/5\|
 \leq(1+h(a_0))\|\Gamma_0F(x_0)\|.
\]
Take $F(u,v)=(u,100v)$, the maximum norm, $x_0=(0,1)$,
$\Gamma_0=\operatorname{diag}(1,1/100)$, $\beta=\eta=1$, and $K=1/10$.
We have $a_0=1/10$ and $y_0=z_0=x_1=0$. The constant derivative
satisfies the Lipschitz bound with this $K$. The inequality would
require $100\leq1+h(1/10)=1689/1600$, which is false. The necessary
bound retains $\Gamma_0$ inside the norm and is proved by first
performing the Newton cancellation. The example does not refute
convergence: the root is reached in the first cycle.

\paragraph{A3. Application of a remainder to a vector.}
In the proof of $(\mathrm{II}_1)$, page 207, a vector appears
multiplied on the right by an integral of operators. In Banach spaces,
the integral must be written as applied to the vector
$-\Gamma_0(9F(y_0)+F(z_0))/5$. The printed identity does not respect
the input and output spaces; correcting it preserves the norm bounds.

\paragraph{Reconstruction of the semilocal step.}
Fix an iteration and set $d=\|\Gamma F(x)\|$, $b\geq\|\Gamma\|$,
$\eta\geq d$, and $a=Kb\eta$. If $Q=(9F(y)+F(z))/5$, the remainders
from $x$ give
\[
 F(y)=E_y,\qquad F(z)=-5F(y)+E_z,\qquad Q=(4F(y)+E_z)/5,
\]
\[
 \|E_y\|\leq Kd^2/2,\quad \|E_z\|\leq K\|z-x\|^2/2,
 \quad \|z-x\|\leq(1+5a/2)d.
\]
Consequently, $\|\Gamma Q\|\leq h(a)d$ and
$\|x^+-x\|\leq(1+h(a))d$. The remainder $E_+$ from $y$ to $x^+$ gives
\[
 \Gamma F(x^+)=\Gamma F(y)-\Gamma Q
 -\Gamma(F'(y)-F'(x))\Gamma Q+\Gamma E_+.
\]
Taking norms,
\[
 \|\Gamma F(x^+)\|\leq[a/2+(1+a)h(a)+ah(a)^2/2]d=g(a)d.
\]
The variation of the derivative and the Banach lemma give
\[
 \|\Gamma^+F'(x)\|\leq f(a),\qquad
 \|\Gamma^+F(x^+)\|\leq f(a)g(a)d,\qquad
 f(a)=[1-a(1+h(a))]^{-1}.
\]
The published functions are recovered without using a fifth-order
assumption. The recurrence $a_{n+1}=a_nf(a_n)^2g(a_n)$ is decreasing
on the interval in Section~\ref{supp:radial}. Formula (8), page 209,
retains a geometric tail that tends to zero. The control of $z_n$
retains $(f(a_0)^3g(a_0)^2)^n$, and the radius includes $5a_0/2$.
The dependence between the domain and the inverses is resolved by
induction over the stages with these same estimates; the limit belongs
to a ball smaller than the one required for the domain.

\paragraph{A4. Uniqueness is indeed proved.}
Theorem 1 states it, and pages 209--210 use the averaged operator
between roots. With the open balls, we obtain
\[
 K\beta\int_0^1\|x_*+t(y_*-x_*)-x_0\|\,dt
 <\frac{K\beta}{2}\left(R\eta+\frac2{K\beta}-R\eta\right)=1.
\]
The Banach lemma applies. Uniqueness is not a conclusion missing from
that work. The remark on its page 210 discusses the scope of the
radii near the threshold.

\paragraph{A5. Order and efficiency indices.}
Theorem 1 does not prove fifth order, although section 4 uses it to
compute indices in finite dimensions. In dimension three, with the
costs of that section, generic order four gives $4^{1/35}<2^{1/17}$;
the claimed superiority over Newton for every dimension greater than
two is not retained. This does not exclude competitiveness in other
problems. Table 2 also computes indices using the total cost, although
the definition uses the cost per iteration. For M5, $640$ operations
per step and four steps produce $2560$ operations: the published value
$1.000629$ corresponds to $5^{1/2560}$, not $5^{1/640}$.
The other rows exhibit the same change of convention.

\paragraph{A6. Reproduction of the eight-node example.}
Gauss--Legendre with $2000$ digits gives
\[
\begin{aligned}
 K&\simeq0.247117984146368,& \beta&\simeq1.596943933144693,\\
 \eta&\simeq0.591064523238323,& a_0&\simeq0.233253900277713.
\end{aligned}
\]
The radii are approximately $1.516495240390509$ and $3.551497385109859$,
in agreement with the published values. Extending the calculation
gives ACOC values $4.7225706005$, $3.8727238969$, $4.0306205065$, and
$3.9925100905$. The residual after four cycles is approximately
$1.2044268509\times10^{-312}$, less than the published tolerance
$10^{-180}$; the number of steps is not refuted. The reproduced first
component is $1.0122394257\ldots$, compared with $1.01222\ldots$ in
their table, a minor discrepancy. The data and the six updates are
included in \path{results/audit/}.

The semilocal core of the 2015 work can be justified with its own
functions. The comparison below uses that core, not the
incorrect order classification or the efficiency tables.

\subsection{Singh et al. (2016): inequalities, radii, and a local counterexample}
\paragraph{B1. H\"older generalization of the auxiliary functions.}
With $q=1$, the functions $t,r,s$ on pages 268--269 of
\cite{SinghEtAl2016} coincide with the 2015 functions $h,f,g$.
For $0<q\leq1$, the estimate of $t$ uses $(1+u)^q\leq1+u^q$ and
does not lose the factor five. The condition
$r(c_0)^{1+q}s(c_0)^q<1$ is consistent with the decrease of $c_n$.

\paragraph{B2. Incorrect power in the semilocal proof.}
Relation (V) gives $(1+t(c_n))d_n$, but (2.17)--(2.20) use
$(1+t(c_n))^qd_n$, which is smaller if $q<1$. The example
\[
\begin{gathered}
 F(u)=u-1-u^2/2000,\quad x_0=0,\quad D=(-1.01,1.01),\\
 \beta=\eta=1,\quad q=1/2,\quad K=c_0=1/700
\end{gathered}
\]
satisfies the semilocal setting. The H\"older condition holds because
$(1/1000)^2(2.02)<(1/700)^2$. The auxiliary contraction also holds
and the correct radius is less than $1.01$; the program checks them
with rational bounds for the radicals. The exact stages give
$x_1=1.000500500625$. If $w=1/1050$,
\[
 t(1/700)=w+w^2+\frac{w(1+5w)}{\sqrt{5250}}
 <w+w^2+\frac{w(1+5w)}{72}<\frac{967}{10^6}.
\]
It follows that
$(1+t(1/700))^{1/2}<1.0004835<1.000500500625=x_1$.
Inequality (2.17) is false. The radius in Theorem 2.1 uses
$1+t(c_0)$ without the power: removing the incorrect exponents
reconstructs the semilocal chain with that radius. The counterexample
refutes the printed inequality; it does not prove the semilocal
conclusion false.

\paragraph{B3. Local uniqueness: a false conclusion of Theorem 3.1.}
The necessary condition $K_0R^q/(1+q)<1$ is equivalent to
$R<((1+q)/K_0)^{1/q}$. Instead, $R<(1+q)/K_0$ is printed.
Consider
\[
 G(u)=u-4u^2,\quad x_*=0,\quad D_0=(-3/10,3/10),\quad
 q=1/2,\quad K_0=9/2,\quad K_1=25/4.
\]
On the domain, $8|u|\leq K_0|u|^{1/2}$ and
$8|u-v|\leq K_1|u-v|^{1/2}$ hold, since
$64(3/10)<(9/2)^2$ and $64(6/10)<(25/4)^2$.
Moreover,
\[
 r_1=\left(\frac{1+q}{K_1+(1+q)K_0}\right)^{1/q}
 =(3/26)^2<7/25.
\]
The local radius $\rho$ used here is less than $r_1$.
The choice $R=7/25$ satisfies $\rho<R<(1+q)/K_0=1/3$ and
$B[0,R]\subset D_0$, but contains the roots $0$ and $1/4$.
Uniqueness in that ball is false. The corrected bound is $1/9$ and
excludes the second root. The refutation specifically concerns local
uniqueness with the printed radius.

\paragraph{B4. Radii of the semilocal examples.}
The expressions for $\eta$ have typographical errors: in Example 2.1,
$16$ appears where $24$ is required in the denominator; in Example 2.2,
a factor $1/4$ is missing from the numerator. Recalculating from the
operators and using the formula in Theorem 2.1 gives:
\begin{table}[htbp]\centering\small
\begin{tabular}{@{}lrrrr@{}}\toprule
Case & Published existence & Recomputed & Published uniqueness & Recomputed\\\midrule
2.1, $q=.85$ & .04630 & .04632668 & 38.25348 & 38.11497803\\
2.1, $q=1$ & .04620 & .04620900 & 21.99790 & 21.95379100\\
2.2, $q=.7$ & .47873 & .49097189 & 6.85383 & 5.89723102\\
2.2, $q=1$ & .44916 & .44915651 & 3.65166 & 3.32162365\\
\bottomrule\end{tabular}
\caption{Multiprecision reproduction of the 2016 radii; these are not
rigorous intervals.}\label{supp:tab:singh-radii}
\end{table}
For $q<1$, the published existence radii are reproduced with the
incorrect power $(1+t(c_0))^q$. The uniqueness radii are reproduced
using
\[
 \eta\left(\frac{1+q}{c_0}-(R\eta)^q\right)^{1/q}
 \quad\text{instead of}\quad
 \eta\left(\frac{1+q}{c_0}-R^q\right)^{1/q}.
\]
The scale $\eta$ is applied twice. This does not exclude justifying
some of the radii by another argument; it indicates that they do not
agree with the theorem's formula. For Example 2.3, the initial vector
yields $\beta\simeq26.5887634801$ and
$\eta\simeq0.000375702982175$.

\paragraph{B5. Local examples and table.}
In Example 3.1, zero is also a root by the definition of $G$;
$x_*=1$ is not the only root in the entire stated interval. This does
not exclude local uniqueness near $1$. In Example 3.2,
$G(u)=2u^{3/2}/3-u$, $x_*=9/4$, and $G'(x_*)=1/2$.
The constant $K_1=1$, with $q=1/2$, does not control all of $[1,3]$,
since $2(\sqrt3-1)>\sqrt2$; the inequality persists at interior points
sufficiently close to the endpoints. The formulas give $r_1=.36$,
not $.361$.
For $q=1/2$, the local radii of Examples 3.1--3.4 are reproduced as
$1.291615219\times10^{-6}$, $0.01206846112$, $0.00194510607$, and
$6.843277886\times10^{-5}$. Table 2 assigns the last three values
to other rows. Fractional powers require an explicit real domain;
varying $q$ in $x^{1+q}$ also changes the equation, not only the theory
applied to a fixed problem.

\paragraph{B6. Recoverable content.}
With $H=1+t(c_0)$ and $C=r(c_0)s(c_0)<1$, the corrected recurrence gives
\[
 \|x_{n+m}-x_n\|\leq H\eta C^n\frac{1-C^m}{1-C}.
\]
The stages are controlled by $H\eta/(1-C)$ and
$5c_0C^n\eta/(1+q)$, which explains the radius in the statement
without the incorrect power. An induction over the stages avoids
circularity between the domain and the inverses. In semilocal
uniqueness, the open balls and $q>0$ allow the strict inequality to be
recovered. The objection concerning merely nondecreasing continuity
in the 2022 work does not automatically transfer to this case.

\subsection{Maroju, Behl, and Motsa (2022): semilocal and local flaws}
\paragraph{C1. Incompatibility of C6 and the recurrence.}
Page 23 of \cite{MarojuBehlMotsa2022} requires
$\varphi\in C[0,1]$, $\varphi\geq1$, and defines
$a_{n+1}=a_nf(a_n)\varphi(c_n)$, $c_n=f(a_n)g(a_n)$.
When $a_n>0$, $f(a_n)>1$, and $c_n\in[0,1]$, we have
$a_{n+1}>a_n$, contrary to Lemma 2.2. If $c_n$ lies outside that
interval, the assumption does not define $\varphi(c_n)$.
The auxiliary functions also use $\varphi(1+h(a))$ and
$\varphi(1+a_0M)$ outside the domain. The condition
$f(a_0)\varphi(a_0)<1$ is also incompatible in the stated regime.
Changing a sign without reconstructing the domain and scaling does
not resolve these dependencies.

\paragraph{C2. Missing factor five and false bound (2.6).}
The estimates on page 24 give
$\|z_0-x_0\|\leq(1+5a_0M)\|y_0-x_0\|$, but (2.6) uses $1+a_0M$.
For $F(u)=u-1-u^2/20$, $x_0=0$, $\beta=\eta=1$, and
$\omega(t)=t/10$, extend $\varphi(s)=\max\{1,s\}$ to $s\geq0$
to isolate that inequality. C6 is retained on $[0,1]$, and
$M=\int_0^1\varphi(t)dt=1$.
The stages are $y_0=1$, $z_0=5/4$, and $a_0=1/10$, but
\[
 \omega(|z_0-x_0|)|z_0-x_0|=\frac5{32}>
 \frac{121}{1000}=\omega(\eta)\varphi(1+a_0M)(1+a_0M)\eta.
\]
The bound is refuted even when only the domain of $\varphi$ is
repaired. No claim is made that the example satisfies the incompatible
contraction conditions in C1.

\paragraph{C3. The function $h$ underestimates the displacement.}
If C6 is repaired as in the Lipschitz case, with $\varphi(s)=s$ and
$M=1/2$, then for $F_a(u)=u-1-au^2/2$ we obtain
\[
 x_1=1+a/2+a^2/2+5a^3/8,\qquad
 h_{\rm imp}(a)=a/2+a^2/10+a^3/40.
\]
Hence $x_1-(1+h_{\rm imp}(a))=a^2(3a+2)/5>0$. At $a=1/10$,
(2.7) would require $1.055625\leq1.051025$.
The choice $\varphi(s)=s$ does not literally satisfy C6: it shows
that this isolated repair does not save the bound.

\paragraph{C4. Residual propagation.}
The definition of $h$ enters (2.11)--(2.12), which are intended to
control $\|\Gamma_1F(x_1)\|$. An estimate compatible with the second
stage would need to retain at least
\[
 h_{\rm corr}(a)=\tfrac45aM+\tfrac15aM\varphi(1+5aM)(1+5aM).
\]
The form $g(a)=aM+(1+a)h(a)+ah(a)M\varphi(h(a))$ may be a reasonable
bound, but it requires a valid $h$ and appropriate scaling conditions.
The missing factor affects the chain that must prove the decay of
the normalized residual.

\paragraph{C5. Mixing recurrences and the Cauchy tail.}
Relation (II) gives factors $f(a_j)g(a_j)$ for $\|\Gamma_nF(x_n)\|$,
not $f(a_0)\varphi(c_0)$, which belongs to another recurrence.
The substitution in (2.15)--(2.17) mixes these quantities.
The sum in (2.17) starts at zero and does not retain the tail that
must tend to zero with its initial index. An appropriate bound would
be of the form $C\eta c_0^m(1-c_0^n)/(1-c_0)$, after proving the
contraction. The published bound does not by itself prove the
Cauchy property.

\paragraph{C6. Containment and domain.}
In (2.19), $1+h(a_0)$ is replaced by $\varphi(1+h(a_0))$ without
an identity that permits it. The last inequality in (2.21) requires
$5Ma_0+R(1-c_0^{n+1})<R$. If $0<c_0<1$ and $Ma_0>0$, the left-hand
side tends to $5Ma_0+R>R$: the estimate does not hold for every index.
The statement does not clearly include containment of the full ball
in the domain, and derivatives are used before containment is
established. The induction needs to be reconstructed, not just a
constant changed.

\paragraph{C7. Mixed regularity assumptions.}
Theorem 2.3 introduces $a_0=K\beta\eta^p$ and uses a H\"older bound
for $F'(x_n)$, although the preceding framework uses
$a_0=\beta\omega(\eta)$ and general $\omega$-continuity.
The consistent bound would be
$\|F'(x_n)\|\leq\|F'(x_0)\|+\omega(R\eta)$.
The identity $F(x_n)=F'(x_n)\Gamma_nF(x_n)$ is correct; the decay
of the second factor remains to be justified with the printed
recurrences.

\paragraph{C8. Semilocal uniqueness.}
Formula (2.22) reaches a bound $\leq1$ and invokes the Banach lemma,
which requires a strict inequality. Nondecreasing continuity of
$\omega$ does not guarantee this merely by using open balls.
Also, $R+r$ appears where the distances are bounded by $R\eta+r$.
A consistent sufficient condition would be
\[
 \beta\int_0^1\omega((1-t)R\eta+tr)\,dt<1,
\]
with the domain inclusions. It is not claimed that uniqueness is
false in every example, but that the printed argument does not
prove it.

\paragraph{C9. Identities and coefficients in the local part.}
Formula (3.14) represents $F(y_0)$ by the segment ending at $x_0$,
whose integral is $F(x_0)-F(x_*)$, not $F(y_0)$; $y_0$ must be used.
There is a similar difficulty in (3.18). In (3.19), the coefficient
$16$ disappears and expressions involving $F'(x_*)F(y_0)$ appear
where the inverse is required to make the spaces compatible.

This does not imply that the printed function $g_3$ cannot be a bound.
Writing $D=1-\omega_0(t)$ and
$A_i=\int_0^1[1+\omega_0(\theta g_i(t)t)]d\theta$, we have
$g_2=g_1+5A_1g_1/D$. The combined stage allows
$g_3^*=g_1+(9A_1g_1+A_2g_2)/(5D)$, whereas the printed formula,
interpreted with the missing parenthesis, is
$g_3=g_2+(A_1g_1+A_2g_2)/(5D)$.
Hence $g_3-g_3^*=17A_1g_1/(5D)\geq0$; a bound can be recovered
by another argument without validating the published identity.
The uniqueness operator on page 31 integrates between $x_*$ and
$2x_*-y_*$, not between the two roots. One must use
$\int_0^1F'(y_*+\theta(x_*-y_*))d\theta$.
The construction of radii does not cover $\omega_0=\omega=0$, when
$g_1=0$ and $g_1-1$ has no positive zero. An infinite-radius
convention or a region where $g_i<1$ is needed.

\paragraph{C10. Local example 3.2 does not satisfy the equation.}
Its operator satisfies $F(0)=-1$, although $x_*=0$ is identified as
a root. The Green kernel vanishes at the endpoints, so every solution
satisfies $x(0)=x(1)=1$ and does not belong to a ball centered at zero
with radius $0.224209$. The real power $x^{3/2}$ is also not defined
on any full open ball around the zero function. In $F'(x)u$, the
factor $u(t)$ is missing under the integral. The printed functions give
\[
 r_0=\left(\frac{\sqrt{137}-3}{4}\right)^2\simeq4.73573753348,
\]
not $3.2000$. The values $r_1\simeq2.63030003815$ and
$r_2\simeq0.448599953403$ are reproduced. For the printed $g_3$,
we obtain $r_3\simeq0.336536461475$; $0.224209$ corresponds to
restoring the coefficient $16$. No adjustment of these radii
resolves $F(0)\neq0$.

\paragraph{C11. Domain of semilocal example 2.5.}
The constants $\beta=40/31$, $\eta=11/62$, and
$a_0\simeq0.137224097318$ are consistent with the initial function
$1$, but no $\varphi$ completing the assumptions is specified.
If $u^{8/5}=|u|^{8/5}$ is interpreted on the whole real space,
\[
 \|F'(x)-F'(y)\|\leq(\|x-y\|^{3/5}+\|x-y\|)/5
\]
fails for the constant functions $x=.01$, $y=-.01$: the two sides
are approximately $0.0257382937792$ and $0.0231270499958$.
Restricting to the positive cone avoids that example, but requires
an open domain and its verification. On a set of diameter $D$,
$\omega(t)=(t^{3/5}+t)/5\leq(1+D^{2/5})t^{3/5}/5$ for $0\leq t\leq D$.
One cannot claim without qualification that every H\"older condition
fails; on a ball separated from zero, the derivative can be Lipschitz.

The particular conclusions admit another check. For
$Tx=1+\mathcal K(|x|^{8/5}+x^2/10)$, $\|\mathcal K\|=1/8$,
on $B[1,0.238277]$, the approximate bounds are
$\|Tx-1\|\leq0.19512846754$ and Lipschitz constant $0.25832039011<1$.
On $B[1,2.54114]$, one constant is $0.51561500993<1$.
\path{code/check_additional_2022.py} verifies with fractions
$\|Tx-1\|<0.196<0.238277$ and constants less than $0.259$ and $0.516$.
The fractional powers are enclosed by comparing integer powers.
This justifies those particular conclusions by a fixed-point argument,
not the general proof or the containment of the method's stages.

\paragraph{C12. Order under weak regularity.}
The local bounds $g_i$ are $O(t)$ in the Lipschitz case and give,
through that route, a quadratic bound, not fifth order. Under
H\"older conditions, they give an exponent $1+q$. Even in dimension
one, $F(u)=u+u|u|$ has a Lipschitz derivative and a simple root, but
from $t>0$ the scheme produces $T(t)=32t^4/5+O(t^5)$.
Its lack of higher regularity matters. Regularity, local order of
a smooth scheme, and order of a semilocal majorization are different
assertions.

\subsection{Assessment and relation to the quantitative comparison}
The defects in the 2015 work are corrected while retaining its
functions and semilocal mechanism, including uniqueness. In the
2016 work, the incorrect power invalidates printed steps of the
proof, but the semilocal radius in the statement admits a
reconstruction; the counterexample in B3 refutes local uniqueness
with the published radius. In the 2022 work, the incompatibilities
of the assumptions and the chain of bounds require a more substantial
reconstruction. In none of these cases does the designation
``fifth order'' replace an independent proof for operators in
Banach spaces.

The comparison in Section~\ref{supp:radial} uses
the corrected semilocal criterion of the 2015 work, not a weakened
version. It distinguishes the radius for all stages, the location
of the limit, and the recurrence-based sum. The advantage is proved
on a common quadratic family and is not inferred from errors found
in other works. The audit is also not an exhaustive priority search.


\section{Applications: derivatives, histories, and reproducibility}\label{supp:applications}
The five problems use the parameters specified below and are formulated
in $C[0,1]$ with the uniform norm. The gamma conditions are checked
for the continuous operators. The following histories come from
solving their discretizations with $N=30$, five cycles, and $1000$
digits. For Chandrasekhar and Bratu, $N=20,40,60$ are also included
with the same precision and number of cycles. The eleven calculations
and profiles are reproduced through
\path{code/reproduce_experiments.py}.
These multiprecision floating-point results are not identified with
the rational checks in Section~\ref{supp:radial}.

\subsection{Continuous verification and parameters of the five problems}


In Problems~1--4 we take $\Omega=C[0,1]$; in the fifth, a positive
domain is specified. By the Lipschitz form of the gamma difference condition established in Section~\ref{supp:radial},
in addition to $\|G(x_0)\|\leq\beta$ and
$B[x_0,R_\gamma]\subset\Omega$, it suffices to verify
\[
 \|G''(x_0)\|\leq2\gamma,\qquad
 \|G''(x)-G''(z)\|\leq L\|x-z\|_\infty,\qquad L\leq6\gamma^2
\]
on $B[x_0,R_\gamma]$, and check
$\alpha=\beta\gamma\leq\alpha_{\rm c}=1547/10000$.
Products, quotients, and exponentials of functions are understood
pointwise. For Problems~1 and~5 we define
\[
 (\mathcal K_0x)(s)=\int_0^1K_0(s,t)x(t)\,dt,\qquad
 K_0(s,t)=\begin{cases}s/(s+t),&s>0,\\0,&s=0.\end{cases}
\]
Although $K_0$ is not continuous at $(0,0)$, $\mathcal K_0$ maps
$C[0,1]$ into itself: continuity for $s>0$ is immediate, and
\[
 |(\mathcal K_0x)(s)|\leq\|x\|_\infty s\log\frac{1+s}{s}
 \longrightarrow0\quad(s\to0^+).
\]
Moreover, $\|\mathcal K_0\|=\log2$.

\medskip
\noindent\textbf{Problem 1. Chandrasekhar's integral equation.}
For the family in \cite{EzquerroHernandezMagrenan2021}, we take
$\omega_0=1/2$, $x_0=1$, and
\[
\begin{aligned}
 F(x)&=x-1-\tfrac{\omega_0}{2}x\mathcal K_0x,\\
 F'(x)h&=h-\tfrac{\omega_0}{2}(h\mathcal K_0x+x\mathcal K_0h),\\
 F''(x)[h,k]&=-\tfrac{\omega_0}{2}(h\mathcal K_0k+k\mathcal K_0h).
\end{aligned}
\]
With $q=\omega_0\log2<1$, we have $\|I-F'(1)\|\leq q$,
$\|F'(1)^{-1}\|\leq(1-q)^{-1}$, $\|F(1)\|_\infty\leq q/2$,
and $\|F''(1)\|\leq q$. We choose
\[
 \beta=\gamma=\frac{q}{2(1-q)}\simeq0.265197109517,\qquad
 \alpha\simeq0.0703295068963<\alpha_{\rm c}.
\]
Thus, $\|G(1)\|\leq\beta$ and $\|G''(1)\|\leq2\gamma$.
Since $G''$ is constant, $L=0$.

\medskip
\noindent\textbf{Problem 2. Polynomial Hammerstein equation.}
Consider the class in \cite{CardenasCastroSierra2022} with
\[
\begin{aligned}
 \relax [F(x)](s)&=x(s)-\sin(s/3)
  -s\int_0^1t^4\bigl(x(t)^3+x(t)^2/2\bigr)\,dt,\\
 [F'(x)h](s)&=h(s)-s\int_0^1t^4(3x(t)^2+x(t))h(t)\,dt,\\
 [F''(x)[h,k]](s)&=-s\int_0^1t^4(6x(t)+1)h(t)k(t)\,dt.
\end{aligned}
\]
For $x_0=0$, $F'(0)=I$, $\|G(0)\|_\infty=\sin(1/3)$, and
$\|G''(0)\|\leq1/5$. Moreover,
$\|G''(x)-G''(z)\|\leq(6/5)\|x-z\|_\infty$. Therefore, we take
\[
 \beta=\sin(1/3),\qquad \gamma=1/\sqrt5,\qquad
 \alpha=\frac{\sin(1/3)}{\sqrt5}\simeq0.146325916783<\alpha_{\rm c},
\]
with $L=6/5=6\gamma^2$ and $\|G''(0)\|<2\gamma$.

The rank-one structure allows a scalar reduction. Let
$\phi(s)=\sin(s/3)$, $\psi(s)=s$, and
$\mathcal L=\{x_a=\phi+a\psi:a\in\mathbb R\}$. Define
\[
 h(a)=a-\int_0^1t^4\left[
 (\phi(t)+a\psi(t))^3+\tfrac12(\phi(t)+a\psi(t))^2\right]dt.
\]
Then $F(x_a)=h(a)\psi$ and $F'(x_a)\psi=h'(a)\psi$.
If $F'(x_a)$ is invertible, $h'(a)\neq0$ and
$F'(x_a)^{-1}F(x_b)=h(b)\psi/h'(a)$ for $x_b\in\mathcal L$.
Since $F'(0)=I$, the first step gives $y_0=\phi$, and also
$z_0,x_1\in\mathcal L$. From $n=1$ onward, the three stages coincide
with the scalar method for $h(a)=0$. Near a simple root, the order is
at least five, and exactly five when the leading coefficient does not
vanish. Quadrature preserves rank one and this reduction.

\medskip
\noindent\textbf{Problem 3. Bratu through Green's function.}
For $x''+\lambda e^x=0$, $x(0)=x(1)=0$, we take $\lambda=2$.
Its integral formulation uses
\[
 (\mathcal Kx)(s)=\int_0^1K(s,t)x(t)\,dt,\qquad
 K(s,t)=\begin{cases}t(1-s),&t\leq s,\\s(1-t),&s\leq t,\end{cases}
\]
with $\|\mathcal K\|=1/8$; see \cite{AhmadArshadUllahMa2023}.
The expressions for $F$ and its derivatives are
\[
 F(x)=x-\lambda\mathcal K(e^x),\quad
 F'(x)h=h-\lambda\mathcal K(e^xh),\quad
 F''(x)[h,k]=-\lambda\mathcal K(e^xhk).
\]
For $x_0=0$, $\|F'(0)^{-1}\|\leq(1-\lambda/8)^{-1}$ and
$\|G(0)\|,\|G''(0)\|\leq\lambda/(8-\lambda)=1/3$.
We choose $\beta=1/3$ and $\gamma=0.4$, so that
$R_\gamma\simeq0.732233047034$ and
\[
 L=\frac{\lambda e^{R_\gamma}}{8-\lambda}
 \simeq0.693239845939<6\gamma^2=0.96,\qquad
 \alpha=\beta\gamma=\frac2{15}<\alpha_{\rm c}.
\]
The Lipschitz bound follows from
$|e^{x(t)}-e^{z(t)}|\leq e^{R_\gamma}|x(t)-z(t)|$ in the ball.

\medskip
\noindent\textbf{Problem 4. Hammerstein with a nonseparable kernel.}
As a specialization of the classes in \cite{EzquerroHernandez2024}, let
$(\mathcal K_{\rm H}x)(s)=\int_0^1e^{st}x(t)\,dt$,
$\lambda=10^{-2}$, and
\[
 [F(x)](s)=x(s)-\sin(s/3)
 -\lambda[\mathcal K_{\rm H}(e^x-1-x)](s).
\]
We have $\|\mathcal K_{\rm H}\|=e-1$,
$F'(x)h=h-\lambda\mathcal K_{\rm H}((e^x-1)h)$, and
$F''(x)[h,k]=-\lambda\mathcal K_{\rm H}(e^xhk)$.
For $x_0=0$, $F'(0)=I$. We take $\beta=\sin(1/3)$ and $\gamma=0.2$,
with $R_\gamma\simeq1.46446609407$. Then
\[
 \|G''(0)\|\leq\lambda(e-1)<0.4=2\gamma,\qquad
 L=\lambda(e-1)e^{R_\gamma}\simeq0.0743196987942<0.24=6\gamma^2,
\]
and $\alpha=\beta\gamma\simeq0.0654389393592<\alpha_{\rm c}$.

\medskip
\noindent\textbf{Problem 5. Reciprocal Hammerstein equation.}
For this class see \cite{AmatBusquierGutierrez2011}. On the open convex
set $\Omega=\{x\in C[0,1]:\min_{t\in[0,1]}x(t)>0\}$, we define
\[
\begin{aligned}
 F(x)&=x-1+\mu\mathcal K_0(1/x),\qquad \mu=0.1,\\
 F'(x)h&=h-\mu\mathcal K_0(h/x^2),\qquad
 F''(x)[h,k]=2\mu\mathcal K_0(hk/x^3).
\end{aligned}
\]
With $x_0=1$, we choose
$\beta=\mu\log2/(1-\mu\log2)\simeq0.0744770755493$ and $\gamma=0.8$.
The Banach lemma gives $\|G(1)\|\leq\beta$ and
$\|G''(1)\|\leq2\beta<2\gamma$.
Since $R_\gamma\simeq0.366116523517<1$, the gamma ball is contained
in $\Omega$. With $\delta=1-R_\gamma$, the inequality
$|x(t)^{-3}-z(t)^{-3}|\leq3\delta^{-4}|x(t)-z(t)|$ gives
\[
 L=\frac{6\mu\log2}{\delta^4(1-\mu\log2)}
 \simeq2.76780915206<3.84=6\gamma^2.
\]
Therefore, $\alpha=\beta\gamma\simeq0.0595816604395<\alpha_{\rm c}$.

\subsection{Discretization, linear solve, and stopping criterion}
Let $(t_j,w_j)_{j=1}^N$ be the Gauss--Legendre nodes and weights on
$[0,1]$, computed with the same precision as the iterations.
The Nystr\"om matrix is defined by $K_{ij}=w_jK(t_i,t_j)$, with the
kernel corresponding to each problem. For the kernel $K_0$, the point
$(0,0)$ is not evaluated, since the Gauss--Legendre nodes are interior.
For the polynomial problem, the coefficients $w_jt_j^4$ are used and
the rank-one structure of the discrete system is preserved exactly.
The Jacobians are evaluated using the preceding analytical formulas,
not finite differences.

In each cycle, the Jacobian $A=F_N'(x_n)$ is factored only once:
$PA=LU$. The three right-hand sides are $F_N(x_n)$, $F_N(y_n)$,
and $9F_N(y_n)+F_N(z_n)$.
For each one, $Lv=Pb$ and $Uu=v$ are solved by reusing $P,L,U$.
The code also records the numerical residual of these linear equations.
This check differs from a proof of backward stability.

We use multiprecision floating-point arithmetic from \texttt{mpmath},
with $1000$ decimal working digits; no double-precision binary
parameters or nodes from a lower-precision rule are introduced.
Exactly five cycles are executed, regardless of the residual.
If $d_n=\|x_n-x_{n-1}\|_\infty$, we compute
\[
 \widehat p_n=\frac{\log(d_n/d_{n-1})}{\log(d_{n-1}/d_{n-2})},\qquad n\geq3.
\]
The quotient is omitted if any of the three increments is zero,
if the denominator is zero, or if an increment does not exceed the
precautionary level $10^{-D+30}\max\{1,\|x_n\|_\infty\}$, where $D$
is the decimal precision. The same comparison marks residuals near
the precision limit; it does not turn them into a rigorous bound
of $10^{-D}$.

\subsection{Profile reconstruction and Bratu reference}
To compare different dimensions, nodal vectors are not subtracted
directly. In Chandrasekhar, we reconstruct
\begin{equation}
 u_N(s)=\left(1-\frac14\sum_{j=1}^Nw_j\frac{s}{s+t_j}x_{N,j}\right)^{-1}
 \quad(s>0),\qquad u_N(0)=1,
 \label{supp:eq:chandra-profile}
\end{equation}
and in Bratu,
\begin{equation}
 u_N(s)=2\sum_{j=1}^Nw_jK(s,t_j)e^{x_{N,j}}.
 \label{supp:eq:bratu-profile}
\end{equation}
These expressions define the Nystr\"om discretizations used in the computations.
The profiles are evaluated with $100$ digits on
$\mathcal G=\{i/400:0\leq i\leq400\}$.

The Bratu reference is
\begin{equation}
 u_*(s)=2\log\frac{\cosh(a/2)}{\cosh(a(s-1/2))},\qquad
 \frac{2a^2}{\cosh(a/2)^2}=2,\quad 0<a<2.
 \label{supp:eq:bratu-exact}
\end{equation}
Indeed, $u_*(0)=u_*(1)=0$,
$u_*''(s)=-2a^2\operatorname{sech}^2(a(s-1/2))$, and
$e^{u_*(s)}=\cosh^2(a/2)\operatorname{sech}^2(a(s-1/2))$;
the relation for $a$ gives $u_*''+2e^{u_*}=0$.
The function $2a^2/\cosh^2(a/2)$ is strictly increasing on $(0,2)$,
since its logarithmic derivative is $2/a-\tanh(a/2)>0$.
The chosen root corresponds to the small branch and is computed
with $100$ digits.

For Chandrasekhar, we report
$\max_{s\in\mathcal G}|u_N(s)-u_{60}(s)|$; for Bratu,
$\max_{s\in\mathcal G}|u_N(s)-u_*(s)|$.
These are maxima on a finite grid. The first uses a numerical
reference, and neither constitutes a bound for the continuous
supremum. The change of slope of the Green kernel along $s=t$ is
retained in the global quadrature used; spectral convergence of
that rule is not assumed.

\subsection{Complete histories for \texorpdfstring{$N=30$}{N=30}}
The tables show the six states $x_0,\ldots,x_5$.
The sign $\dagger$ indicates that the residual reached the
precautionary precision level and is not interpreted as an exact
bound. The ACOC of the fifth update uses the preceding increments,
which do remain above that level. The vectors of all stages and
the values without presentation rounding are included in the
JSON and CSV files.

\begin{table}[htbp]\centering\small
\begin{tabular}{@{}rlll@{}}\toprule
$n$ & $\|F_N(x_n)\|_\infty$ & $d_n$ & $\widehat p_n$\\\midrule
0 & $1.732117\times10^{-1}$ & -- & --\\
1 & $1.583097\times10^{-6}$ & $2.51117\times10^{-1}$ & --\\
2 & $7.980421\times10^{-27}$ & $2.358968\times10^{-6}$ & --\\
3 & $1.832861\times10^{-111}$ & $9.793127\times10^{-27}$ & $4.054341706$\\
4 & $1.486176\times10^{-450}$ & $1.991172\times10^{-111}$ & $4.15526647$\\
5 & $\dagger$ & $1.80824\times10^{-450}$ & $4.003242371$\\
\bottomrule\end{tabular}
\caption{Chandrasekhar: $N=30$, five cycles, and $1000$ digits.}\label{supp:tab:history-chandrasekhar}
\end{table}
\begin{table}[htbp]\centering\small
\begin{tabular}{@{}rlll@{}}\toprule
$n$ & $\|F_N(x_n)\|_\infty$ & $d_n$ & $\widehat p_n$\\\midrule
0 & $3.267054\times10^{-1}$ & -- & --\\
1 & $8.26069\times10^{-6}$ & $3.400329\times10^{-1}$ & --\\
2 & $3.491167\times10^{-28}$ & $9.104126\times10^{-6}$ & --\\
3 & $4.706249\times10^{-140}$ & $3.847615\times10^{-28}$ & $4.893408432$\\
4 & $2.095062\times10^{-699}$ & $5.186758\times10^{-140}$ & $5.000002792$\\
5 & $\dagger$ & $2.308968\times10^{-699}$ & $5.0$\\
\bottomrule\end{tabular}
\caption{Polynomial Hammerstein: $N=30$, five cycles, and $1000$ digits.}\label{supp:tab:history-polinomico}
\end{table}
\begin{table}[htbp]\centering\small
\begin{tabular}{@{}rlll@{}}\toprule
$n$ & $\|F_N(x_n)\|_\infty$ & $d_n$ & $\widehat p_n$\\\midrule
0 & $2.497788\times10^{-1}$ & -- & --\\
1 & $2.523014\times10^{-6}$ & $3.288667\times10^{-1}$ & --\\
2 & $2.896231\times10^{-26}$ & $3.553876\times10^{-6}$ & --\\
3 & $1.165631\times10^{-105}$ & $3.38339\times10^{-26}$ & $4.031427609$\\
4 & $1.329004\times10^{-423}$ & $1.620031\times10^{-105}$ & $3.961762207$\\
5 & $\dagger$ & $1.549674\times10^{-423}$ & $4.009328914$\\
\bottomrule\end{tabular}
\caption{Bratu--Green: $N=30$, five cycles, and $1000$ digits.}\label{supp:tab:history-bratu}
\end{table}
\begin{table}[htbp]\centering\small
\begin{tabular}{@{}rlll@{}}\toprule
$n$ & $\|F_N(x_n)\|_\infty$ & $d_n$ & $\widehat p_n$\\\midrule
0 & $3.267054\times10^{-1}$ & -- & --\\
1 & $4.771053\times10^{-10}$ & $3.271316\times10^{-1}$ & --\\
2 & $1.276932\times10^{-45}$ & $4.787416\times10^{-10}$ & --\\
3 & $1.467877\times10^{-187}$ & $1.279543\times10^{-45}$ & $4.026550214$\\
4 & $1.155496\times10^{-755}$ & $1.472923\times10^{-187}$ & $3.990068027$\\
5 & $\dagger$ & $1.157858\times10^{-755}$ & $4.002458988$\\
\bottomrule\end{tabular}
\caption{Nonseparable Hammerstein: $N=30$, five cycles, and $1000$ digits.}\label{supp:tab:history-no_separable}
\end{table}
\begin{table}[htbp]\centering\small
\begin{tabular}{@{}rlll@{}}\toprule
$n$ & $\|F_N(x_n)\|_\infty$ & $d_n$ & $\widehat p_n$\\\midrule
0 & $6.928469\times10^{-2}$ & -- & --\\
1 & $4.112001\times10^{-10}$ & $7.287908\times10^{-2}$ & --\\
2 & $2.875395\times10^{-44}$ & $4.298189\times10^{-10}$ & --\\
3 & $9.334016\times10^{-179}$ & $2.963562\times10^{-44}$ & $4.151191393$\\
4 & $7.602065\times10^{-719}$ & $9.75626\times10^{-179}$ & $3.936672652$\\
5 & $\dagger$ & $7.835002\times10^{-719}$ & $4.016099651$\\
\bottomrule\end{tabular}
\caption{Reciprocal Hammerstein: $N=30$, five cycles, and $1000$ digits.}\label{supp:tab:history-reciproco}
\end{table}
\subsection{Refinement of Chandrasekhar and Bratu}
The following table is computed using the reconstructions
\eqref{supp:eq:chandra-profile}--\eqref{supp:eq:bratu-profile}.
All iterations use $1000$ digits; only the evaluation of the
profiles and the Bratu reference uses $100$ digits.
\begin{table}[htbp]\centering\small
\begin{tabular}{@{}lrrl@{}}\toprule
Problem & $N$ & $\widehat p_5$ & Difference on $\mathcal G$\\\midrule
Chandrasekhar & 20 & 4.003222946 & $5.697734303\times10^{-5}$\\
Chandrasekhar & 30 & 4.0032423709 & $8.08036804\times10^{-6}$\\
Chandrasekhar & 40 & 4.0029442845 & $1.100911601\times10^{-6}$\\
Chandrasekhar & 60 & 4.0031870353 & reference\\
Bratu--Green & 20 & 4.0098824137 & $1.770558072\times10^{-3}$\\
Bratu--Green & 30 & 4.0093289144 & $7.862010152\times10^{-4}$\\
Bratu--Green & 40 & 4.00912517 & $4.431748795\times10^{-4}$\\
Bratu--Green & 60 & 4.008982179 & $1.990539703\times10^{-4}$\\
\bottomrule\end{tabular}
\caption{Maxima on the grid of $401$ points. The reference is $N=60$
for Chandrasekhar and the explicit small branch for Bratu.}
\label{supp:tab:refinement}
\end{table}
The rounded values in this table agree with the stored histories and profiles. The observed stability of the ACOC
and reduction of the profile differences do not replace a rigorous
analysis of the Nystr\"om error.

\subsection{Files, environment, and scope}
The directory \path{results/experiments/} contains eleven
subdirectories, one per problem, number of nodes, and precision.
Each includes \path{result.json} (parameters, environment, and
summary), \path{history.csv}, \path{nodes_weights_solution.csv},
and \path{complete_iterates.json}. The latter also contains the
intermediate stages of the five cycles. The aggregate files are
\path{five_problems_N30.csv}, \path{refinement.csv},
\path{profiles_grid401.csv}, and \path{Bratu_reference.json}.

The included execution uses Python 3.13.5, \texttt{mpmath} 1.3.0 and
its Python backend, on Linux x86\_64. Each record stores the platform,
the full version, and the SHA-256 hash of the program. No performance
conclusions are drawn from the times, since the hardware conditions
and system load were not controlled.

The programs allow the observations for the discrete system to be
reproduced, but do not certify rounding errors, quadrature errors,
or errors in the continuous supremum. Such validation would require
additional bounds or interval arithmetic. The conclusions for the
exact iteration in $C[0,1]$ follow from the verified assumptions and
the semilocal criterion.


\section{Methodological distinction from previous perturbed majorizations}
\label{supp:perturbations}
We document the comparison with Argyros \cite{Argyros1996} and
Ferreira--Svaiter \cite{FerreiraSvaiter2012}. The formulas from
those sources are distinguished from the algebraic correspondences
we establish for the scheme studied here. The comparison is
limited to the identified constructions and results; the letters
used to describe each previous work belong to that context and
are not automatically identified with our controls.

\subsection{The framework of Argyros includes exact corrections}
On pages 120 and 132 of \cite{Argyros1996}, a Newton stage and a
Newton-like stage, respectively, are followed by a correction.
Calling this correction $d_n$, the structure is
\[
 y_n=x_n-A(x_n)^{-1}F(x_n),\qquad x_{n+1}=y_n-d_n.
\]
The author includes $A(x_n)=F'(x_n)$ and exact choices of $d_n$.
It is not appropriate to describe that entire framework as an
analysis of external computational errors.
Section II, page 121, incorporates an additive term $\alpha_n$
into the scalar remainder:
\[
 t_{n+1}=s_n+I_n,\qquad
 h_{n+1}=\frac L2(t_{n+1}-s_n)^2+\alpha_n,\qquad
 s_{n+1}=t_{n+1}+\frac{h_{n+1}}{1-Lt_{n+1}}.
\]
Accumulated bounds for $\sum\alpha_i$ and $\sum I_i$ are required.
In the application in its Section VI, page 136, it provides
\[
 \alpha_n=\frac N6(s_n-t_n)^3+M(s_n-t_n)(t_{n+1}-s_n),
\]
with derivative control constants. This is an explicit precedent
for additive terms constructed with scalar increments for an exact
correction. Its Remark 1(a), page 123, also uses radial functions
instead of a fixed Lipschitz constant.

\paragraph{Correspondence with our scheme.}
This observation is a deduction from the comparison, not an assertion
in \cite{Argyros1996}. For our method, let
\[
 Q_n=\frac{9F(y_n)+F(z_n)}5,\qquad
 \Gamma_n=F'(x_n)^{-1},\qquad d_n=\Gamma_nQ_n.
\]
Then $x_{n+1}=y_n-d_n$ formally has the preceding structure.
This does not verify the assumptions on majorants, sums, and domains
of Argyros's theorem, nor does it automatically locate the additional
point $z_n$.

The relation also appears in the remainder. Set $p_n=y_n-x_n$,
$q_n=z_n-y_n$, $r_n=x_{n+1}-y_n$, $D_n=F'(y_n)-F'(x_n)$,
\[
 \mathcal B_n(d)=F(y_n+d)-F(y_n)-F'(y_n)d,\qquad w_n=r_n-q_n/5.
\]
The term in condition (4) of \cite{Argyros1996}, evaluated at
$(x_n,y_n,x_{n+1})$, contains
\[
 F(y_n)-F(x_n)-F'(x_n)(y_n-x_n)+F'(y_n)r_n
 =F(y_n)+F'(y_n)r_n,
\]
because the Newton step cancels $F(x_n)+F'(x_n)p_n$.
For the scheme studied, the residual identity gives
\[
 F(y_n)+F'(y_n)r_n=D_nw_n-\tfrac15\mathcal B_n(q_n).
\]
Its estimate can use the cancellations in \eqref{supp:eq:residual-identities}.
The distinction is not the prior absence of an additive remainder,
but its specific control together with the location of the three
stages and the comparison with the fixed scalar trajectory.

\subsection{The Ferreira--Svaiter slack is not the realized residual}
In sections 4 and 5 of \cite{FerreiraSvaiter2012}, a step $s$ from
$x$ with tolerance $\theta\geq0$ satisfies
\[
 \|F'(x_0)^{-1}(F(x)+F'(x)s)\|
 \leq\theta\|F'(x_0)^{-1}F(x)\|.
\]
The analysis introduces
\[
 K(t,\epsilon)=\{x:\|x-x_0\|\leq t,\quad
 \|F'(x_0)^{-1}F(x)\|\leq f(t)+\epsilon\}
\]
and, in formula (4.3) of that source,
\begin{equation}
 t^+=t-(1+\theta)\frac{f(t)+\epsilon}{f'(t)},\qquad
 \epsilon^+=\epsilon+2\theta(f(t)+\epsilon).
 \label{supp:eq:FS-update}
\end{equation}
The relation $N_\theta(K(t,\epsilon))\subset K(t^+,\epsilon^+)$
is the invariance mechanism of its Proposition 4.5. It is not merely
a similarity in notation with $g(t_n)+\widehat\Delta_n$: there is
a coupled construction of a radius and a residual slack.

The variable $\epsilon$ belongs to the proof; it is not the residual
$F(x)+F'(x)s$ or its norm. Its increment depends on the allowed
tolerance and the preceding residual bound. With $\theta=0$ and
$\epsilon_0=0$, the recurrence gives $\epsilon_n=0$. This does not
extend to an initial value $\epsilon_0>0$, nor does it mean that
choosing an exact step under a positive tolerance forces the slack
to vanish.

Its Theorem 5.1 establishes that $\epsilon_n$ is nondecreasing and
bounded. For $\theta>0$ and a positive initial residual bound,
$\epsilon_1>\epsilon_0\geq0$, so $\sum\epsilon_n$ does not converge.
The increments can, however, be summable. The quantity tending to
zero is $f(t_n)+\epsilon_n$, not necessarily $\epsilon_n$ separately.
Here, instead, $\sum\widehat\Delta_n<\infty$ and
$\sum\widehat E_n<\infty$ are proved.
The difference describes the role of the variables; it does not
by itself establish superiority between the theorems.

In its application to analytic functions, page 361, it uses
\[
 f(t)=\frac{t}{1-\gamma t}-2t+b
 =b-t+\frac{\gamma t^2}{1-\gamma t},
\]
the same rational form as our majorant. The assumptions and the
methods controlled are different, so their criteria are not
identified, but this rational form is not claimed as a new feature
of the present construction.

\subsection{Rewriting as inexact Newton: scope and limits}
For the total step $s_n^{\rm tot}=x_{n+1}-x_n=p_n+r_n$, we have
$F'(x_n)p_n=-F(x_n)$ and $F'(x_n)r_n=-Q_n$. Hence
\begin{equation}
 F(x_n)+F'(x_n)s_n^{\rm tot}=-Q_n.
 \label{supp:eq:inexact-rewriting}
\end{equation}
The exact execution of a multipoint method does not prevent its
total step from being represented as an inexact correction relative
to another algorithm, in this case one-step Newton.
A uniform admissible relative tolerance does not follow from
\eqref{supp:eq:inexact-rewriting}. One would also have to check
\[
 \|F'(x_0)^{-1}Q_n\|\leq\theta\|F'(x_0)^{-1}F(x_n)\|
\]
at all indices, with a tolerance satisfying the assumptions of the
result invoked. Two separate upper bounds for the numerator and
denominator do not prove that relative comparison. The existence
and location of $y_n,z_n$ needed to construct $Q_n$ also require
their own argument. Neither the impossibility nor the general
validity of that semilocal reduction is claimed here.

\subsection{Meaning and logical role of our defects}
The hatted quantities are constructed from $\alpha$ through the
polynomial first cycle and the rational recurrences in
Section~\ref{supp:radial}. The sequences $t_n,s_n,v_n$ are obtained
by applying the same scalar scheme to $g$. They are compared through
\[
 \widehat\Delta_n=\frac{[\widehat\varepsilon_n-\gamma g(t_n)]_+}{\gamma},
 \qquad
 \widehat E_n=\frac{[\widehat A_n+\widehat C_n-
 \gamma(t_{n+1}-t_n)]_+}{\gamma}.
\]
For these fixed controls, they are the smallest nonnegative slacks,
since
\begin{equation}
 \begin{aligned}
 \gamma(g(t_n)+\widehat\Delta_n)
 &=\max\{\gamma g(t_n),\widehat\varepsilon_n\},\\
 \gamma(t_{n+1}-t_n+\widehat E_n)
 &=\max\{\gamma(t_{n+1}-t_n),\widehat A_n+\widehat C_n\}.
 \end{aligned}\label{supp:eq:defect-identities}
\end{equation}
The actual norms can be smaller than their bounds: a positive defect
does not prove that the actual residual violates exact majorization.
Nor is it a realized physical or computational perturbation.

Summability is not obtained merely by writing
\eqref{supp:eq:defect-identities}. First, closure of the controls and
$\sum_{n\geq0}(\widehat A_n+\widehat C_n)=\widehat\tau_\infty<\infty$
are proved. The inequalities
$\widehat\varepsilon_n\leq\widehat A_n$,
$\gamma\widehat\Delta_n\leq\widehat\varepsilon_n$, and
$\gamma\widehat E_n\leq\widehat A_n+\widehat C_n$
then give the convergent series, also from $n=0$.

The error bounds have different roles. Summing the second identity
in \eqref{supp:eq:defect-identities},
\begin{equation}
 \begin{aligned}
 \gamma\left(t_*-t_n+\sum_{k=n}^{\infty}\widehat E_k\right)
 ={}&\widehat\tau_\infty-\widehat\tau_n\\
 &+\sum_{k=n}^{\infty}
 [\gamma(t_{k+1}-t_k)-(\widehat A_k+\widehat C_k)]_+.
 \end{aligned}\label{supp:eq:two-bounds}
\end{equation}
The radius bound does not exceed the one written with defects.
The latter retains the reference to the scalar trajectory; the
former uses the constructed controls directly. Both depend only
on the initial data and are a priori. The estimate using a verified
residual $\|G(x_n)\|$ is a posteriori. The quantitative improvement
is not attributed to the mere use of the positive part.

\subsection{Contribution supported by these comparisons}
The previous works contain ingredients of a perturbed majorization
theory. We do not claim that every additive term is new or that the
literature lacks deterministic controls for exact algorithms.
Reading two articles is also not sufficient to prove priority
throughout a whole class.

The construction combines a fixed scalar reference of the three-step
method, residual bounds that retain the cancellations, stage radii,
a specific estimate for the first cycle, and finite tail control.
Its result is the criterion $\beta\gamma\leq1547/10000$, the error
bounds, and the quantitative comparison in Section~\ref{supp:radial}.
Applying a more general formalism might require precisely these
estimates. Until such an application is proved, neither complete
equivalence nor irreducibility is claimed.

\begin{thebibliography}{99}
\interlinepenalty=10000 % Keep each bibliographic entry on one page.
\bibitem{CorderoSemilocal2015} A. Cordero, M. A. Hern\'andez-Ver\'on, N. Romero and J. R. Torregrosa, Semilocal convergence by using recurrence relations for a fifth-order method in Banach spaces, \textit{Journal of Computational and Applied Mathematics}, 273, 205--213, 2015. \url{https://doi.org/10.1016/j.cam.2014.06.008}.

\bibitem{SinghEtAl2016} S. Singh, D. K. Gupta, E. Mart\'inez and J. L. Hueso, Semilocal and local convergence of a fifth order iteration with Fr\'echet derivative satisfying H\"older condition, \textit{Applied Mathematics and Computation}, 276, 266--277, 2016. \url{https://doi.org/10.1016/j.amc.2015.11.062}.

\bibitem{MarojuBehlMotsa2022} P. Maroju, R. Behl and S. S. Motsa, Semilocal and local convergence of a three step fifth order iterative methods under general continuity condition in Banach spaces, \textit{Thai Journal of Mathematics}, 20(1), 21--33, 2022. \url{https://thaijmath.com/index.php/thaijmath/article/view/1307}.

\bibitem{Argyros1996} I. K. Argyros, Results on Newton methods: Part I. A unified approach for constructing perturbed Newton-like methods in Banach space and their applications, \textit{Applied Mathematics and Computation}, 74(2--3), 119--141, 1996. \url{https://doi.org/10.1016/0096-3003(95)00090-9}.

\bibitem{FerreiraSvaiter2012} O. P. Ferreira and B. F. Svaiter, A robust Kantorovich's theorem on the inexact Newton method with relative residual error tolerance, \textit{Journal of Complexity}, 28(3), 346--363, 2012. \url{https://doi.org/10.1016/j.jco.2012.02.002}.

\bibitem{AhmadArshadUllahMa2023} J. Ahmad, M. Arshad, K. Ullah and Z. Ma, Numerical solution of Bratu's boundary value problem based on Green's function and a novel iterative scheme, \textit{Boundary Value Problems}, 2023, Article 102, 2023. \url{https://doi.org/10.1186/s13661-023-01791-6}.

\bibitem{AmatBusquierGutierrez2011} S. Amat, S. Busquier and J. M. Guti\'errez, Third-order iterative methods with applications to Hammerstein equations: A unified approach, \textit{Journal of Computational and Applied Mathematics}, 235(9), 2936--2943, 2011. \url{https://doi.org/10.1016/j.cam.2010.12.011}.

\bibitem{CardenasCastroSierra2022} E. C\'ardenas, R. Castro and W. Sierra, Convergence analysis for the King--Werner method under gamma-conditions, \textit{Journal of Applied Mathematics and Computing}, 68, 4605--4620, 2022. \url{https://doi.org/10.1007/s12190-022-01720-3}.

\bibitem{EzquerroHernandez2024} J. A. Ezquerro and M. A. Hern\'andez-Ver\'on, Location, separation and approximation of solutions of nonlinear Hammerstein-type integral equations, \textit{Applied Numerical Mathematics}, 198, 1--10, 2024. \url{https://doi.org/10.1016/j.apnum.2023.12.010}.

\bibitem{EzquerroHernandezMagrenan2021} J. A. Ezquerro, M. A. Hern\'andez-Ver\'on and A. A. Magre\~n\'an, On an efficient modification of the Chebyshev method, \textit{Computational and Mathematical Methods}, 3(6), e1187, 2021. \url{https://doi.org/10.1002/cmm4.1187}.
\end{thebibliography}
\end{document}
